%% TRACK A COMPRESSED: sections 1-6 and 9 of the Witten two-track restructure. Merge pass pending.

\providecommand{\cZ}{\mathcal{Z}}

% ====================================================================
\section*{1.\quad The bar construction on algebraic curves}
% ====================================================================

\medskip

\noindent
The classical bar construction of a graded algebra lives on the
cofree tensor coalgebra $T^c(s^{-1}\bar A)$; its differential encodes
the multiplication. For chiral algebras on a curve~$X$ three features
distinguish the theory. The multiplication is mediated by a
logarithmic propagator $d\log(z_i-z_j)$, and the bar differential
becomes an integral transform on the Fulton--MacPherson
compactification of configuration spaces. The bar coalgebra carries
two coproducts: a non-cocommutative deconcatenation on ordered
configurations $\operatorname{Conf}_n^{<}(\mathbb R)$, producing the
primitive $E_1$-coalgebra $\barB^{\mathrm{ord}}(\cA)=T^c(s^{-1}\bar\cA)$,
and a cocommutative coshuffle on the Ran space, producing the symmetric
factorization coalgebra $\barB^{\Sigma}(\cA)=\operatorname{Sym}^c(s^{-1}\bar\cA)$
as its $\Sigma_n$-coinvariant shadow. On curves of genus $g\ge 1$ the
propagator deforms via the prime form, the bar differential acquires
curvature $\kappa(\cA)\cdot\omega_g$ from the Hodge bundle, and the
universal Maurer--Cartan element~$\Theta_{\cA}$ of the modular
convolution algebra controls the genus tower.

The ordered bar is the primitive object of this monograph. The
symmetric bar, and with it every statement about the modular
characteristic~$\kappa$ and the shadow obstruction tower, is obtained
as the image under a canonical averaging map
$\mathrm{av}\colon\barB^{\mathrm{ord}}(\cA)\to\barB^{\Sigma}(\cA)$
described in \S1.6.

\subsection*{1.1.\enspace The $1$-form}

Let $X$ be a smooth algebraic curve over~$\mathbf C$. On the
configuration space $C_2(X)=X\times X\setminus\Delta$, define the
logarithmic $1$-form
\begin{equation}\label{eq:preface-eta}
\eta_{12}\;=\;d\log(z_1-z_2)\;=\;\frac{dz_1-dz_2}{z_1-z_2}.
\end{equation}
This form has a simple pole along~$\Delta$ with residue~$1$. The
residue depends on no coordinate, no metric, no level. A double
pole~$dz/(z-w)^2$ transforms under coordinate change and has no
coordinate-independent residue; a regular $1$-form extracts nothing
from the collision. The logarithmic derivative is the unique $1$-form
on $C_2(X)$ with a first-order pole along~$\Delta$ and conformally
invariant residue.

On $C_3(X)$, the three pullbacks $\eta_{12}, \eta_{23}, \eta_{31}$
satisfy the \emph{Arnold relation}
\begin{equation}\label{eq:preface-arnold}
\eta_{12}\wedge\eta_{23}+\eta_{23}\wedge\eta_{31}+\eta_{31}\wedge\eta_{12}\;=\;0,
\end{equation}
the differential-form image of the partial-fractions identity
$\tfrac{1}{(z_1-z_2)(z_2-z_3)}+\tfrac{1}{(z_2-z_3)(z_3-z_1)}+\tfrac{1}{(z_3-z_1)(z_1-z_2)}=0$.
Arnold proved that~\eqref{eq:preface-arnold}, replicated across all
triples $\{i,j,k\}\subset\{1,\ldots,n\}$, generates the complete ideal
of relations in $H^*(\mathrm{Conf}_n(\mathbf C),\mathbf C)$.

Under $z\mapsto w(z)$ the propagator transforms as
$\eta_{ij}\mapsto\eta_{ij}+d\log\frac{w(z_i)-w(z_j)}{z_i-z_j}$;
near the diagonal this reduces to $d\log w'(z_i)$, the standard
$\mathrm{Diff}(X)$ cocycle. The propagator is a BRST ghost for
diffeomorphisms.

\subsection*{1.2.\enspace Chiral algebras on a curve}

Write $\mathcal D_X$ for the sheaf of differential operators on~$X$.
A \emph{chiral algebra}~$\cA$ is a $\mathcal D_X$-module equipped with
a \emph{chiral bracket}
\[
\mu^{\mathrm{ch}}\colon j_*j^*(\cA\boxtimes\cA)\longrightarrow\Delta_!(\cA),
\]
where $j\colon X\times X\setminus\Delta\hookrightarrow X\times X$ is
the inclusion of the complement of the diagonal and $\Delta_!\cA$ is
the $!$-pushforward along the diagonal; concretely $\Delta_!\cA\cong
\cA\otimes\delta_\Delta$. The source $j_*j^*(\cA\boxtimes\cA)$ consists
of sections of the external product with poles of arbitrary order
along~$\Delta$; the chiral bracket extracts the singular part and
projects it onto the diagonal.

In local coordinates the chiral bracket encodes the \emph{operator
product expansion}
\[
a(z)\,b(w)\;\sim\;\sum_{n\ge 0}\frac{a_{(n)}b(w)}{(z-w)^{n+1}}.
\]
The $n$-th products $a_{(n)}b\in\cA$ are extracted by
$a_{(n)}b=\operatorname{Res}_{z=w}[a(z)b(w)(z-w)^n]$. The zeroth
product is the Lie-type bracket; the first product carries the
symmetric/metric data; higher products ($n\ge 2$) are determined by
the $\mathcal D_X$-module structure. The \emph{Borcherds identity}
couples all pole orders into a single constraint,
\[
\sum_{j\ge 0}\binom{m}{j}(a_{(n+j)}b)_{(m+k-j)}c
=\sum_{j\ge 0}(-1)^j\binom{n}{j}\bigl(a_{(m+n-j)}(b_{(k+j)}c)-(-1)^n b_{(n+k-j)}(a_{(m+j)}c)\bigr),
\]
specialising at $n=0$ to Jacobi, at $m=0$ to the commutator formula,
at $n=1$, $m=k=0$ to the derivation property of~$a_{(0)}$.

\subsection*{1.3.\enspace The bar complex on~$X$}

Place sections $a_1,\dots,a_n\in\cA$ at distinct points
$z_1,\dots,z_n\in X$ and assign to each ordered pair $(i,j)$ the
form $\eta_{ij}=d\log(z_i-z_j)$. For a spanning tree $T$ of the
complete graph on $\{1,\ldots,n\}$ with edge set $E(T)$ form
\[
\omega_T(a_1,\ldots,a_n)=a_1(z_1)\otimes\cdots\otimes a_n(z_n)\otimes\bigwedge_{(i,j)\in E(T)}\eta_{ij},
\]
a section of $\Omega^{n-1}_{\log}$ of degree $|E(T)|=n-1$ on
$\mathrm{Conf}_n(X)$.

The Fulton--MacPherson compactification
$\overline{\mathrm{FM}}_n(X)$ resolves $\mathrm{Conf}_n(X)$ by
iterated real oriented blowup along all diagonal strata: a blowup,
not a complement. Points of
$\overline{\mathrm{FM}}_n(X)\setminus\mathrm{Conf}_n(X)$ record
limiting direction and relative scale as points collide; boundary
codimension-one faces are indexed by subsets
$S\subseteq\{1,\ldots,n\}$ with $|S|\ge 2$.

The \emph{ordered chiral bar complex} on~$X$ (the primitive object):
\[
\barB^{\mathrm{ord}}_X(\cA)
\;=\;
\bigoplus_{n\ge 1}(s^{-1}\bar\cA)^{\otimes n}\otimes
\Gamma\bigl(\overline{\mathrm{FM}}_n(X),\Omega^{n-1}_{\log}\bigr).
\]
Here $\bar\cA=\ker(\varepsilon\colon\cA\to\omega_X)$ is the
augmentation ideal (using $\bar\cA$ rather than $\cA$ is essential:), and $s^{-1}$ is desuspension with
$|s^{-1}a|=|a|-1$ in cohomological convention ($|d|=+1$). The
deconcatenation coproduct
\[
\Delta[s^{-1}a_1|\cdots|s^{-1}a_n]
=\sum_{i=0}^n[s^{-1}a_1|\cdots|s^{-1}a_i]\otimes[s^{-1}a_{i+1}|\cdots|s^{-1}a_n]
\]
splits an ordered sequence at every position. The differential
decomposes as
\[
d_{\barB}=d_{\mathrm{int}}+d_{\mathrm{res}}+d_{\mathrm{dR}}.
\]
The three components: $d_{\mathrm{int}}$ is the internal differential
of~$\cA$; $d_{\mathrm{res}}$ extracts OPE residues along
codimension-one boundary strata of $\overline{\mathrm{FM}}_n(X)$ at
the pole order dictated by $\bigwedge\eta_{ij}$; $d_{\mathrm{dR}}$ is
the de Rham differential on the logarithmic forms.

\subsection*{1.4.\enspace $d^2=0$ and Arnold+Borcherds}

Expand $(d_{\mathrm{int}}+d_{\mathrm{res}}+d_{\mathrm{dR}})^2$. The
diagonal terms: $d_{\mathrm{int}}^2=0$, $d_{\mathrm{dR}}^2=0$; the
cross-terms $[d_{\mathrm{int}},d_{\mathrm{dR}}]=0$ by Leibniz. The
remaining five terms pair into three cancellation mechanisms:

\emph{Disjoint supports}: when $\{i,j\}\cap\{k,l\}=\emptyset$ the
boundary faces are transverse, and the residue operators commute.

\emph{Shared index}: when $\{i,j\}\cap\{k,l\}=\{j=k\}$ the triple
collision $z_i=z_j=z_l$ lies in a codimension-two stratum. The
form contribution $\eta_{ij}\wedge\eta_{jl}$ paired with its cyclic
partners vanishes by the Arnold relation; the corresponding
algebraic contractions
$(a_{i(m)}a_j)_{(n)}a_l+(a_{j(m)}a_l)_{(n)}a_i+(a_{l(m)}a_i)_{(n)}a_j$
vanish by Borcherds. Arnold kills the form; Borcherds kills the
algebra; both are consumed.

\emph{Same pair}: $d_{\mathrm{res}}^{(ij)}\,d_{\mathrm{res}}^{(ij)}=0$
by the desuspension sign rule.

Result: $d_{\barB}^2=0$ iff the $n$-th products satisfy the Borcherds
identity. The bar construction accepts exactly chiral algebras.

All of this is genus~$0$. At genus $g\ge 1$ the propagator acquires
quasi-periodicity and $d_{\mathrm{fib}}^2\neq 0$: the curvature
$d_{\mathrm{fib}}^2=\kappa(\cA)\cdot\omega_g$ produces a single scalar
governing the genus expansion. Genus~$0$ is uncurved algebra;
genus~$\ge 1$ is curved geometry.

\subsection*{1.5.\enspace The two dualities of the bar complex}

The cobar functor $\Omega(\barB(\cA))$ is the categorical logarithm
inverted on the \emph{Koszul locus}: the full subcategory on which
\[
\varepsilon\colon\Omega(\barB(\cA))\xrightarrow{\;\sim\;}\cA
\]
is a quasi-isomorphism. This is bar-cobar \emph{inversion}, recovering
$\cA$ from its bar complex (Theorem~B). The Koszulness question, which
controls the scope of this inversion, is the central organizing
property of the programme and is characterised by twelve equivalent
conditions in~\S9.

A second, linear operation on bar cohomology produces the \emph{Koszul
dual algebra}~$\cA^!$. First, take bar cohomology
$\cA^{\mathrm i}:=H^*(B(\cA))$, a graded coalgebra, and apply graded
linear duality: $\cA^!:=(\cA^{\mathrm i})^\vee$. The cobar recovers
the original; linear duality of bar cohomology produces a different
algebra except at self-dual fixed points (e.g.\ $\mathrm{Vir}_{c=13}$).

A third, geometric operation acts on the bar coalgebra. Verdier
duality on $\mathrm{Ran}(X)$ sends the factorization coalgebra
$\barB_X^\Sigma(\cA)$ to the \emph{homotopy Koszul dual}~$\cA^!_\infty$,
a factorization \emph{algebra} carrying full chain-level data. The
cohomology of $\cA^!_\infty$ recovers $\cA^!$ on the Koszul locus;
this identification is the content of Theorem~A, the Verdier
intertwining theorem.

On $X=\mathbb A^1$, classical Koszul duality (Priddy,
Beilinson--Ginzburg--Soergel, Loday--Vallette) embeds: the affine
line deformation-retracts to a point (but the retraction is
additional data, not a canonical identification), the FM chain
cooperad becomes quasi-isomorphic to the associative cooperad, and
the bar construction collapses to the classical bar complex of
the underlying Lie algebra. On~$\mathbb{P}^1$, the global
topology (Arnold relations, Fulton--MacPherson compactification)
has no counterpart over a bare point.

\subsection*{1.6.\enspace $E_1$ primacy and the averaging map}

The bar complex of a chiral algebra carries two operadic colours, each
with its own Koszul duality.

\medskip
\noindent\textbf{The ordered bar (primitive).}\enspace
The Stasheff tensor bar $\barB^{\mathrm{ord}}(\cA)=T^c(s^{-1}\bar\cA)$
with deconcatenation coproduct is an $E_1$-coalgebra: coassociative,
not cocommutative. It is built on ordered configurations
$\operatorname{Conf}_n^{<}(\mathbb R)$ parametrising positions of line
operators on an interval; the linear ordering is retained, with no
$\Sigma_n$-quotient. Its Koszul dual $\cA^!_{\mathrm{line}}$ governs
the category of line operators: for affine Kac--Moody,
$\cA^!_{\mathrm{line}}$ is the Yangian $Y(\mathfrak g)$.

\medskip
\noindent\textbf{The symmetric bar (coinvariant shadow).}\enspace
The symmetric bar $\barB^{\Sigma}(\cA)=\operatorname{Sym}^c(s^{-1}\bar\cA)$
is the $\Sigma_n$-coinvariant quotient of the ordered bar, with the
coshuffle coproduct: $2^n$ bipartitions of an unordered set rather
than $n+1$ consecutive splittings. This is an $E_\infty$-coalgebra,
built on $\overline{\operatorname{FM}}_n(\mathbb C)$ with
$\Sigma_n$-coinvariants: the Beilinson--Drinfeld factorization
coalgebra on $\mathrm{Ran}(X)$. Its Koszul dual $\cA^!_{\mathrm{ch}}$
is the Francis--Gaitsgory chiral Koszul dual via Verdier duality on
$\operatorname{Ran}(X)$.

\medskip
\noindent\textbf{The averaging map and $R$-matrix descent.}\enspace
The averaging map
\[
\mathrm{av}\colon\barB^{\mathrm{ord}}(\cA)\longrightarrow\barB^{\Sigma}(\cA)
\]
takes $\Sigma_n$-coinvariants. For a chiral algebra with OPE poles,
this descent is not naive: it uses the spectral $R$-matrix $R(z)$ as
twisting datum. The $R$-matrix is the degree-$2$ collision residue of
the ordered Maurer--Cartan element,
\[
r(z)\;=\;\operatorname{Res}^{\mathrm{coll}}_{0,2}\bigl(\Theta^{E_1}_{\cA}\bigr),
\]
a matrix-valued rational function encoding the full binary OPE data.
Applying $\mathrm{av}$ at degree~$2$ collapses the profile to a single
scalar,
\begin{equation}\label{eq:av-organizing}
\mathrm{av}\bigl(r(z)\bigr)\;=\;\kappa(\cA),
\end{equation}
the modular characteristic of Theorem~D. At degree~$3$ the KZ
associator projects to the cubic shadow~$\mathfrak C$; at degree~$4$
the Yangian higher coproduct projects to the quartic resonance
class~$\mathfrak Q$. Each degree-$n$ component of the ordered
Maurer--Cartan element projects under $\mathrm{av}$ to the degree-$n$
shadow of the modular obstruction tower $\Theta_{\cA}$.

\medskip
\noindent\textbf{The five theorems as averaged invariants.}\enspace
The main theorems A--D and H are the invariants of the ordered
Maurer--Cartan element $\Theta^{E_1}_{\cA}$ that survive averaging.
The modular characteristic~$\kappa$, the shadow obstruction tower, the
complementarity pairing, and the chiral Hochschild cohomology are all
$\Sigma_n$-coinvariant projections of the richer $E_1$-data. The
ordered story is the primitive; the modular story is its shadow. The
$R$-matrix is derived from the local OPE for vertex algebras
($E_\infty$-chiral) and is independent input for nonlocal quantum
vertex algebras ($E_1$-chiral, in the sense of Etingof--Kazhdan).

\subsection*{1.7.\enspace The Heisenberg atom}

The Heisenberg chiral algebra $\mathcal H_k$ is generated by a single
field~$\alpha$ of conformal weight~$1$ with OPE
$\alpha(z)\alpha(w)\sim k/(z-w)^2$. Only $\alpha_{(1)}\alpha=k\cdot\mathbf 1$
is nonzero; the Borcherds identity is satisfied trivially.

The ordered bar complex at tensor degree~$n$ is
$(s^{-1}\alpha)^{\otimes n}\otimes\Gamma(\overline{\mathrm{FM}}_n(X),\Omega^{n-1}_{\log})$.
At degree~$2$, $d_{\mathrm{res}}(s^{-1}\alpha\otimes s^{-1}\alpha\otimes\eta_{12})=k\cdot\mathbf 1$.
At degree~$n\ge 3$, every residue at a codimension-one collision
stratum extracts $\alpha_{(1)}\alpha=k\cdot\mathbf 1$; the remaining
$n-2$ factors are~$\alpha$, and $\mathbf 1_{(p)}\alpha=0$ for $p\ge 0$
kills further contractions. All structure is captured at tensor
degree~$2$: the bar complex is quadratic, and $\mathcal H_k$ is
chirally Koszul. The modular characteristic is $\kappa(\mathcal H_k)=k$.

The Koszul dual is the curved symmetric chiral algebra
\[
\mathcal H_k^!=\bigl(\operatorname{Sym}^{\mathrm{ch}}(V^*),\;m_0=-k\,\omega\bigr),
\]
a degree-$2$ curvature satisfying the curved $A_\infty$ relation
$m_1^2(a)=[m_0,a]$. The curvature is the bar-dual of the OPE coefficient.

\medskip
\noindent\textbf{The $E_1$ face of Heisenberg.}\enspace
The ordered bar carries the $R$-matrix
$r(z)=k/z$, a scalar with a simple pole: the second-order OPE pole
$k/(z-w)^2$ drops by one order under the $d\log$ kernel. Applying \eqref{eq:av-organizing},
$\mathrm{av}(r(z))=k=\kappa(\mathcal H_k)$. For Heisenberg the
$R$-matrix \emph{is} the modular characteristic: averaging loses
nothing, because $r(z)$ is already $\Sigma_2$-invariant. For the
affine Kac--Moody algebra $\widehat{\mathfrak g}_k$, the $R$-matrix is
the level-full Casimir form
\[
r(z)\;=\;k\,\Omega/z,\qquad\Omega=\sum_a J^a\otimes J_a
\]. Averaging collapses the Casimir
to its trace:
\[
\mathrm{av}(k\Omega/z)\;=\;\frac{(k+h^\vee)\dim\mathfrak g}{2h^\vee}\;=\;\kappa(\widehat{\mathfrak g}_k).
\]
The Casimir tensor structure, invisible to~$\kappa$, is the data
that builds the Yangian.

\subsection*{1.8.\enspace The Kac--Moody OPE}

The affine Kac--Moody algebra $\widehat{\mathfrak g}_k$ has OPE
\[
J^a(z)J^b(w)\sim\frac{f^{ab}{}_{c}J^c(w)}{z-w}+\frac{k\kappa^{ab}}{(z-w)^2},
\]
with $f^{ab}{}_c$ the structure constants and $\kappa^{ab}$ the Killing
form normalised so long roots have squared length~$2$. The first pole
carries the Lie bracket $J^a_{(0)}J^b=f^{ab}{}_cJ^c$; the second
carries the invariant form $J^a_{(1)}J^b=k\kappa^{ab}\cdot\mathbf 1$.
The bar differential decomposes as
\[
\mu^{\mathrm{ch}}=\underbrace{f^{ab}{}_c J^c\cdot\eta_{12}}_{\text{Lie bracket}}+\underbrace{k\kappa^{ab}\cdot\mathbf 1\cdot d\eta_{12}}_{\text{inner product}}.
\]
The Lie bracket contributes the tree-level bar differential
(Chevalley--Eilenberg combinatorics); the inner product contributes
the curvature~$\kappa$. The coexistence of both poles, absent in the
pure second-order Heisenberg, makes $\widehat{\mathfrak g}_k$ the
first nontrivial example.

\begin{remark}[The zero-dimensional shadow]\label{rem:preface-zero-dim}
When $X=\operatorname{Spec}\Bbbk$ is a point, configuration spaces
collapse, logarithmic forms vanish, and the chiral bar complex reduces
to the classical bar construction of an associative algebra~$A$. A
quadratic algebra $A=T(V)/(R)$ has quadratic dual
$A^!=T(V^*)/(R^\perp)$; when $A$ is Koszul, bar cohomology concentrates
in diagonal degrees and $H^*(B(A))\cong(A^!)^\vee$. The operadic
instances $\mathrm{Com}^!=\mathrm{Lie}$ and $\mathrm{Sym}^!=\Lambda$
recover the Chevalley--Eilenberg complex as a restricted bar
construction. When Koszulness fails, the cobar-bar resolution
$\Omega(B(A))\xrightarrow{\sim}A$ is the universal replacement.
\end{remark}

Everything above takes place at genus~$0$: the propagator is globally
defined on $\mathbb P^1$, and the Arnold relation holds without
correction. On a curve of genus~$g\ge 1$, $z_1-z_2$ has no global
meaning; the curve is compact, and any meromorphic function with a
simple zero must also have a pole. The Arnold relation acquires a
defect, the bar differential no longer squares to zero, and the
defect is controlled by a single scalar.


% ====================================================================
\section*{2.\quad Curvature and the genus tower}
% ====================================================================

\subsection*{2.1.\enspace The genus-$g$ propagator}

At genus $g\ge 1$, the replacement for $z_1-z_2$ is the \emph{prime
form} $E(z_1,z_2)$: a section of $K^{-1/2}\boxtimes K^{-1/2}$
vanishing to first order along the diagonal with no other zeros. The
bundle is $K^{-1/2}\boxtimes K^{-1/2}$; the prime form is a
$(-\tfrac12,-\tfrac12)$-differential. In a local coordinate with
diagonal $z_1=z_2$,
\[
E(z_1,z_2)=(z_1-z_2)\bigl(1+O(z_1-z_2)\bigr)\cdot(dz_1)^{-1/2}(dz_2)^{-1/2}.
\]
Analytic continuation around the $B$-cycle $B_j$ in~$z_1$ gives
\[
E(z_1+B_j,z_2)=-\exp\bigl(-\pi i\Omega_{jj}-2\pi i\textstyle\int_{z_2}^{z_1}\omega_j\bigr)\cdot E(z_1,z_2),
\]
where $\Omega=(\Omega_{jk})$ is the period matrix and
$\omega_1,\ldots,\omega_g$ the holomorphic differentials normalised by
$\oint_{A_j}\omega_k=\delta_{jk}$. The $A$-cycle monodromy is trivial.

The logarithmic derivative
$\omega(z_1,z_2):=d_{z_1}\log E(z_1,z_2)$ is the fundamental
meromorphic differential of the second kind. The \emph{genus-$g$
propagator}:
\[
\eta^{(g)}_{12}=d\log E(z_1,z_2)+\pi\sum_{j,k=1}^g(\operatorname{Im}\Omega)^{-1}_{jk}\operatorname{Im}\!\Bigl(\int_{z_2}^{z_1}\omega_j\Bigr)\omega_k(z_1).
\]
The period correction $(\operatorname{Im}\Omega)^{-1}_{jk}$ compensates
$B$-cycle monodromy, rendering $\eta^{(g)}_{12}$ single-valued; this
is a full $g\times g$ matrix. The price:
$\eta^{(g)}_{12}$ is no longer holomorphic. The propagator couples to
the Hodge bundle $\mathbb E\to\overline{\mathcal M}_g$ through
$\Omega$.

\subsection*{2.2.\enspace Curvature}

The Arnold relation at genus~$g$ acquires a defect. The period
corrections produce a residual $(1,1)$-form, and the fibrewise bar
differential satisfies
\begin{equation}\label{eq:pref-curvature}
d_{\mathrm{fib}}^{2}\;=\;\kappa(\cA)\cdot\omega_g,
\end{equation}
where $\omega_g$ is the Arakelov $(1,1)$-form on the fibre of the
universal curve $\pi\colon\mathcal C_g\to\overline{\mathcal M}_g$, and
$\kappa(\cA)\in\Bbbk$ is the \emph{modular characteristic}, determined
entirely by the leading OPE singularity.

For Heisenberg, every residue extracts $\alpha_{(1)}\alpha=k$ and
$\kappa(\mathcal H_k)=k$. For Kac--Moody, the curvature receives the
trace over the adjoint representation:
\[
\kappa(\widehat{\mathfrak g}_k)=\frac{(k+h^\vee)\dim\mathfrak g}{2h^\vee}.
\]
At the critical level $k=-h^\vee$ the shifted level vanishes and the
bar complex is uncurved. The Sugawara construction is undefined (the
central charge $c=k\dim\mathfrak g/(k+h^\vee)$ has a pole) but the
curvature parameter vanishes.

For Virasoro at central charge~$c$, the OPE is
$T(z)T(w)\sim\tfrac{c/2}{(z-w)^4}+\tfrac{2T(w)}{(z-w)^2}+\tfrac{\partial T(w)}{z-w}$.
The bar propagator absorbs one pole order, so the collision
residue is $r(z)=(c/2)/z^3+2T/z$ (cubic, not quartic:).
The curvature computation gives
\[
\kappa(\mathrm{Vir}_c)=c/2.
\]
The Koszul dual is $\mathrm{Vir}_{26-c}$, with
$\kappa(\mathrm{Vir}_{26-c})=(26-c)/2$; the two sum to~$13$, not zero.
At $c=13$, Virasoro is self-dual under Koszul duality.

\subsection*{2.3.\enspace Period integrals restore nilpotence}

The curvature~\eqref{eq:pref-curvature} obstructs $d_{\mathrm{fib}}^2=0$
but not the total $D_g^2=0$. Integration against $A$-cycle periods
defines a connection on the family of bar complexes parametrised by
$\overline{\mathcal M}_g$. The corrected total differential
\[
D_g=d_{\mathrm{fib}}+\nabla^{\mathrm{GM}}
\]
incorporates the Gauss--Manin connection on the Hodge bundle, and
$D_g^2=0$: the $(1,1)$-curvature of the fibrewise differential is
absorbed by the $(0,1)$-part of the base connection. The fibrewise
class $\kappa\cdot\omega_g\in H^{1,1}(\mathcal C_g)$ is identified
under the period map $\overline{\mathcal M}_g\to\mathcal A_g$ with
the $(0,1)$-curvature of $\nabla^{\mathrm{GM}}$ on $R^1\pi_*\mathbb C$,
and the two curvatures cancel.

\subsection*{2.4.\enspace The genus tower (uniform weight)}

The bar complex traces a family of curved cochain complexes over
$\overline{\mathcal M}_g$. For \emph{uniform-weight} algebras
(single-generator, or multi-generator with all generators of the same
conformal weight), the obstruction to flatness is the characteristic
class
\begin{equation}\label{eq:obs-uniform}
\operatorname{obs}_g(\cA)=\kappa(\cA)\cdot\lambda_g\in H^*(\overline{\mathcal M}_g),\qquad\lambda_g=c_g(\mathbb E),
\qquad \end{equation}
where $\mathbb E=\pi_*\omega_\pi$ is the Hodge bundle: the vector
bundle on $\overline{\mathcal M}_g$ whose fibre over $[C]$ is the
$g$-dimensional space $H^0(C,\omega_C)$. At genus~$1$ the
factorization holds unconditionally for all families. For
\emph{multi-weight} families at genus~$g\ge 2$ the scalar formula
\emph{fails}: a cross-channel correction from mixed-propagator
boundary graphs appears,
\begin{equation}\label{eq:obs-multiweight}
\operatorname{obs}_g(\cA)=\kappa(\cA)\cdot\lambda_g+\delta F_g^{\mathrm{cross}},\qquad\text{(all-weight)}
\end{equation}
(Theorem~\ref{thm:multi-weight-genus-expansion}). Every
appearance of $\operatorname{obs}_g$, $F_g$, or $\lambda_g$ in this
monograph carries explicit tagging: \textsc{(uniform-weight)} or
\textsc{(all-weight, with cross-channel correction)}.

Grothendieck--Riemann--Roch applied to the universal curve
$\pi\colon\mathcal C_g\to\overline{\mathcal M}_g$ computes the total
obstruction: the Chern character of $R\pi_*\omega_\pi$ involves the
Todd class of the relative tangent, the $\hat A$-genus. The
generating function in the uniform-weight lane:
\[
\sum_{g\ge 1}\operatorname{obs}_g(\cA)\cdot x^{2g}=\kappa(\cA)\cdot\bigl(\hat A(ix)-1\bigr).
\]
Since $\hat A(ix)=(x/2)/\sin(x/2)=1+x^2/24+7x^4/5760+\cdots$, the first
Faber--Pandharipande coefficients are
\[
F_1=\tfrac{\kappa}{24},\qquad F_2=\tfrac{7\kappa}{5760},\qquad F_3=\tfrac{31\kappa}{967680},
\qquad \]
with the general term
$\kappa\cdot\tfrac{2^{2g-1}-1}{2^{2g-1}}\cdot\tfrac{|B_{2g}|}{(2g)!}$.
Each of these is the scalar-lane projection; for multi-weight
families at $g\ge 2$, add $\delta F_g^{\mathrm{cross}}$. Verify
$F_1=\kappa/24$ as sanity check$, not $1/(2\pi)$).

The scalar $\kappa(\cA)$ is only the leading term. The bar complex
carries higher-degree data: a universal Maurer--Cartan element
$\Theta_\cA$ whose projection to degree~$r$ yields a shadow invariant
$S_r(\cA)$. On each primary deformation line~$L$, the shadow
generating function $H(t)=\sum_{r\ge 2}rS_rt^r$ satisfies the
algebraic relation
\begin{equation}\label{eq:preface-riccati-hook}
H(t)^2=t^4\,Q_L(t),\qquad Q_L(t)=4\kappa^2+12\kappa\alpha\,t+(9\alpha^2+16\kappa S_4)\,t^2,
\end{equation}
a quadratic polynomial in three genus-$0$ invariants
(Theorem~\ref{thm:riccati-algebraicity}). The infinite tower is
algebraic of degree~$2$, governed by the discriminant
$\Delta=8\kappa S_4$.

The shadow obstruction tower is not merely formal. Its projections
yield: the Pixton ideal of the modular curve, the integrable
hierarchies of Gelfand--Dickey type, the Hitchin spectral curve of
the shadow connection, the Baxter $Q$-operator, and the
Bekenstein--Hawking entropy formula of three-dimensional gravity
(Chapter~\ref{ch:twisted-holography-quantum-gravity}). These are
projections of one object, not separate results.

\subsection*{2.5.\enspace Complementarity}

Verdier duality on the Ran space splits the center local system into
two halves, one controlled by $\cA$ and one by $\cA^!$, carrying a
shifted-symplectic structure:
\[
R\Gamma(\overline{\mathcal M}_g,\mathcal Z_\cA)\simeq\mathbf Q_g(\cA)\oplus\mathbf Q_g(\cA^!).
\]
What $\cA$ sees as deformation, $\cA^!$ sees as obstruction. The
natural pairing
\[
\mathbf Q_g(\cA)\otimes\mathbf Q_g(\cA^!)\to R\Gamma(\overline{\mathcal M}_g,\omega_{\overline{\mathcal M}_g})\simeq\Bbbk[-(3g-3)]
\]
is a $(-(3g-3))$-shifted symplectic form, and the two summands are
Lagrangian: isotropic of half the total dimension. The scalar sum
rule is family-dependent (never universal):
\[
\kappa(\cA)+\kappa(\cA^!)=\begin{cases}0&\text{(Kac--Moody, free field, lattice)},\\ 13&\text{(Virasoro)},\\ 250/3&(\mathcal W_3),\\ K_N\cdot(H_N-1)&(\mathcal W_N).\end{cases}
\]
The constant $\kappa+\kappa'$ is the Koszul conductor of the family;
it is zero for KM and free fields, nonzero for $\mathcal W$-algebras.
Qualify the duality statement: ``self-dual at $c=13$ under Koszul
duality,'' never unqualified.

\subsection*{2.6.\enspace Feigin--Frenkel duality}

The Sugawara field
$T(z)=\tfrac{1}{2(k+h^\vee)}\sum_a{:}J^a(z)J^a(z){:}$
has central charge $c=k\dim\mathfrak g/(k+h^\vee)$. At the critical
level $k=-h^\vee$ the denominator vanishes: $T(z)$ is undefined, not
``$c$ diverges.'' The chiral algebra acquires the Feigin--Frenkel
center $\mathfrak z(\widehat{\mathfrak g}_{-h^\vee})$, isomorphic to
functions on the space of ${}^L\mathfrak g$-opers on the disc.

The Feigin--Frenkel involution $k\leftrightarrow-k-2h^\vee$ is Koszul
duality on configuration spaces: it interchanges $\cA$ and $\cA^!$
while preserving the bar-cobar adjunction. At the critical level, the
Koszul dual is the Langlands-dual critical-level
algebra~$\widehat{{}^L\mathfrak g}_{-{}^Lh^\vee}$, not a renormalised
copy of the original.

\subsection*{2.7.\enspace The total space}

The universal configuration fibration
$\pi_n\colon\overline{\mathcal C}_n(\mathcal C_g)\to\overline{\mathcal M}_g$
packages the bar complex into a single global object. Its boundary
divisors carry two distinct geometric roles. \emph{Collision
divisors} (vertical): the fibrewise boundary of
$\overline{\mathrm{FM}}_n(X_s)$, producing the bar differential
$d_{\mathrm{res}}$ (OPE residues). \emph{Degeneration divisors}
(horizontal): the boundary of $\overline{\mathcal M}_g$ itself (nodal
curves), producing the period corrections and Gauss--Manin terms.

Algebraic boundary (Borcherds) meets geometric boundary (period map
and clutching). Their interaction is the curvature~$\kappa$: the
fibrewise bar differential fails to be nilpotent precisely because
the vertical and horizontal boundaries do not commute.


% ====================================================================
\section*{3.\quad Modular homotopy theory}
% ====================================================================

The genus tower demands an organising framework. The fibrewise
differential, the Gauss--Manin correction, and the clutching morphisms
at boundary strata satisfy compatibility relations dictated by the
stratification of $\overline{\cM}_{g,n}$. The governing algebraic
structure is the \emph{modular operad}: an operad whose operations are
indexed by connected stable graphs of arbitrary genus.

The \emph{modular homotopy type} of $\cA$ is the formal moduli problem
classifying the inverse system of bar complexes,
\[
\mathcal M^{\mathrm{mod}}_\cA(R):=\operatorname{MC}(\mathfrak g^{\mathrm{mod}}_\cA\widehat\otimes_{\mathfrak m}R)/{\sim}
\]
for local Artin rings~$R$. Here $\mathfrak g^{\mathrm{mod}}_\cA$ is
the \emph{modular convolution algebra}: a complete filtered cyclic
$L_\infty$-algebra. Three inputs enter: a homotopy chiral algebra on
the curve, a cooperadic chain complex on logarithmic configuration
spaces, and a convolution $L_\infty$-structure mediating between them.
The bar complex, the Koszul dual, the genus tower, the
complementarity splitting, and the shadow obstruction tower are all
representations of $\mathcal M^{\mathrm{mod}}_\cA$.

The primitive story lives on ordered configurations. The $E_1$
convolution algebra $\mathfrak g^{\mathrm{mod},E_1}_\cA$, built on
ribbon graphs from $F\!\Ass$, is the primary object; the modular
convolution algebra $\gAmod$ is its image under the av map from
ordered to symmetric. Every construction in this section begins with
the $E_1$ structure and descends.

\subsection*{3.1.\enspace The modular operad}

At genus~$0$, the combinatorics is governed by trees; at higher
genus, by stable graphs with genus labels at their vertices. A
\emph{modular operad} (Getzler--Kapranov~\cite{GeK98}) has operations
indexed by connected stable graphs. A \emph{connected stable graph}
$\Gamma$ has a label $g(v)\in\mathbb Z_{\ge 0}$ at each vertex
satisfying $2g(v)-2+\operatorname{val}(v)>0$ and total genus
$g(\Gamma)=b_1(\Gamma)+\sum_v g(v)$; trees are the genus-$0$ case.
The structure maps are \emph{separating gluing}
$\xi_{\mathrm{sep}}\colon\cO(g_1,n_1{+}1)\otimes\cO(g_2,n_2{+}1)\to\cO(g_1{+}g_2,n_1{+}n_2)$
and \emph{non-separating gluing}
$\xi_{\mathrm{ns}}\colon\cO(g,n{+}2)\to\cO(g{+}1,n)$.

A stable graph records a degeneration of a smooth curve into a nodal
curve: vertices are irreducible components, edges are nodes, genus
labels are component genera, half-edges are marked points. The set
$\mathsf{Gr}^{\mathrm{st}}_{g,n}$ parametrises the boundary
stratification of $\overline{\cM}_{g,n}$. Every codimension-one
boundary divisor is a one-edge degeneration; every codimension-two
stratum is the intersection of exactly two codimension-one divisors,
appearing with opposite orientations.

The \emph{Feynman transform} replaces trees by stable graphs,
\begin{equation}\label{eq:preface-feynman-sum}
F\cO(g,n)=\bigoplus_{\Gamma\in\mathsf{Gr}^{\mathrm{st}}_{g,n}}\frac{\hbar^{g(\Gamma)}}{|\operatorname{Aut}(\Gamma)|}\bigotimes_{v\in V(\Gamma)}\cO\bigl(g(v),\operatorname{In}(v)\bigr).
\end{equation}

\subsection*{3.2.\enspace The Feynman transform differential}

The differential contracts one edge at a time,
\[
d_{F\cO}(\gamma_\Gamma)=\sum_{e\in E(\Gamma)}(-1)^{\epsilon(e)}\xi_e(\gamma_{\Gamma/e}),
\]
with the sign from the position of $e$ in the ordering on $E(\Gamma)$
recorded by $\det(E(\Gamma))$. The identity $d^2=0$ is the
codimension-$2$ cancellation: each pair of edges contracts in either
order, and the two terms cancel by associativity of~$\cO$. At
genus~$0$, $F\cO$ reduces to the operadic bar construction.

When $\cO=\Com$, each vertex factor in~\eqref{eq:preface-feynman-sum}
is one-dimensional, and the Feynman transform becomes the stable
graph complex
\[
F\Com(g,n)=\bigoplus_{\Gamma\in\mathsf{Gr}^{\mathrm{st}}_{g,n}}\frac{\hbar^{g(\Gamma)}}{|\operatorname{Aut}(\Gamma)|}\det(E(\Gamma)).
\]
When $\cO=\Ass$, the Feynman transform $F\!\Ass$ sums over ribbon
graphs (fatgraphs); this is the $E_1$-primitive input for the ordered
convolution algebra.

\subsection*{3.3.\enspace The convolution ladder}

At each level of generality, the bar-type construction is a
convolution dg Lie algebra whose Maurer--Cartan elements classify
algebraic structures:
\begin{equation}\label{eq:preface-convolution-ladder}
\begin{array}{r@{\;:\quad}c@{\qquad}l}
\text{algebra} & \Hom(B(\Ass),\End_A) & A_\infty\text{-structures on }A\\[3pt]
\text{operad} & \Hom_\Sigma(B(\cP),\End_V) & \cP_\infty\text{-structures on }V\\[3pt]
\text{modular} & \Hom_\Sigma(F\cO,\End_{\barB(\cA)}) & \text{consistent genus towers}
\end{array}
\end{equation}
The first line is Hochschild; the second is the Ginzburg--Kapranov
deformation complex; the third governs this book. The Lie bracket is
the convolution bracket
\[
[f,g]=\mu\circ(f\otimes g)\circ\Delta-(-1)^{|f||g|}\mu\circ(g\otimes f)\circ\Delta,
\]
with $\Delta$ the cooperadic decomposition and $\mu$ the operadic
composition.

The \emph{primitive $E_1$ convolution}
\[
\mathfrak g^{\mathrm{mod},E_1}_\cA:=\Hom_\Sigma\bigl(F\!\Ass,\End_{\barB^{E_1}(\cA)}\bigr)
\]
is the ordered, ribbon-graph input: the algebraic form of the 't~Hooft
$1/N$ expansion. Setting $\hbar=1/N^2$, the MC element
$\Theta^{E_1}_\cA\in\operatorname{MC}(\mathfrak g^{\mathrm{mod},E_1}_\cA)$
is the genus expansion $\sum_g N^{2-2g}F_g(\lambda)$ summed over
ribbon graphs weighted by the genus of their embedding surface.
Planar ($g=0$) diagrams dominate at large~$N$; the shadow obstruction
tower $\Theta^{E_1,\le r}$ is the truncation to $r$-point 't~Hooft
interactions.

The \emph{modular convolution algebra} is the av-image,
\begin{equation}\label{eq:preface-gmod}
\gAmod:=\Hom_\Sigma\bigl(F\Com,\End_{\barB(\cA)}\bigr).
\end{equation}
The coinvariant map $F\!\Ass\to F\Com$ projects
$\Theta^{E_1}\mapsto\Theta$: colour-ordered amplitudes to full
amplitudes. On each primary line, the shadow tower is algebraic of
degree~$2$: the entire infinite sequence is the Taylor expansion of
$\sqrt{Q_L}$ for a single quadratic $Q_L$, the shadow metric
(\S\ref{sec:shadow-metric}).

\medskip
\noindent\textbf{Remark} (Homotopy types in three stages). The three
levels of the convolution ladder mirror three levels of homotopy
theory. \emph{Over a field} (the bar complex of a graded algebra):
$B(A)=(T^c(s^{-1}\bar A),d_B)$ determines
the homotopy type of an associative algebra; formality corresponds to
intrinsic formality; Massey products measure the obstruction.
\emph{Over a manifold}: the Sullivan minimal model $(\Lambda V,d)$
determines the rational homotopy type; K\"ahler manifolds are formal.
The Heisenberg algebra is the vertex-algebraic analogue: its shadow
obstruction tower terminates at degree~$2$, the Gaussian archetype.
\emph{Over a curve at all genera}: the $A_\infty$-products $m_k$ of
the first stage become the modular $L_\infty$-brackets $\ell_k^{(g)}$;
formality becomes the Gaussian archetype; the shadow obstruction tower
is the modular Massey product tower. The shadow depth classification
(\S6) makes this analogy precise: the four classes G, L, C, M
correspond to four levels of modular formality.

\medskip
\noindent\textbf{Koszulness as central question.}\enspace
Within this formality classification, the Koszulness property is the
central organizing question of the programme. A chiral algebra is
\emph{chirally Koszul} iff its bar cohomology concentrates in
diagonal degrees and the cobar-bar counit
$\Omega(\barB(\cA))\to\cA$ is a quasi-isomorphism. Twelve equivalent
conditions (ten unconditional, one conditional on perfectness, one
one-directional) characterise this locus. The Koszulness programme
(\S9) proves the core equivalences from PBW concentration, bar-cobar
inversion, Hochschild polynomial growth, and Koszul-complex
acyclicity. All standard families are Koszul; shadow depth classifies
the complexity of Koszul algebras, not Koszulness itself.

\subsection*{3.4.\enspace Homotopy chiral algebras}

Fix an affine cover $\mathfrak U=\{U_\alpha\}$ of~$X$. The \v Cech
totalisation
\[
A^{\mathrm{ch}}_\infty:=\check C^\bullet(\mathfrak U;\cA)=\bigoplus_{p\ge 0}\prod_{\alpha_0<\cdots<\alpha_p}\cA(U_{\alpha_0}\cap\cdots\cap U_{\alpha_p})
\]
carries a $\mathrm{Ch}_\infty$-algebra structure (Malikov--Schechtman):
the chiral bracket restricted to overlaps satisfies Jacobi up to a
coherent system of higher homotopies, the \emph{secondary Borcherds
operations}
\[
F_n\colon\operatorname{Lie}(n)\otimes\cY(n)\to P_n(\{A^{\mathrm{ch}}_\infty\},A^{\mathrm{ch}}_\infty),\qquad n\ge 2,
\]
encoding derived OPE data at all pole orders. $F_2$ is the chiral
bracket; $F_3$ is the Jacobiator homotopy arising from
$\overline{\mathcal M}_{0,4}\cong\mathbb P^1$, whose boundary has
three points corresponding to OPE compositions
$(a_{(m)}b)_{(n)}c$ and cyclic permutations. Higher operations $F_n$
are determined inductively by $C_*(\overline{\mathcal M}_{0,n+1})$;
the transferred operations are the associahedral $A_\infty$-corrections.

Different affine covers yield quasi-isomorphic $\mathrm{Ch}_\infty$-algebras.
The strict chiral algebra $\cA$ is the $H^0$ of this structure: the
primitive local object is the homotopy chiral algebra, and the strict
algebra appears after rectification. Vallette rectification: the
homotopy category of $\mathrm{Ch}_\infty$-algebras is equivalent to
the homotopy category of strict chiral algebras, via the bar-cobar
adjunction.

\subsection*{3.5.\enspace Logarithmic Fulton--MacPherson spaces}

The classical FM compactification records collisions on a smooth
variety. Its logarithmic extension (Mok) records collisions along a
divisor and provides the normal-crossings boundary structure on which
$D^2=0$ depends.

For a simple-normal-crossing pair $(X,D)$, Mok constructs
$\mathrm{FM}_n(X|D)$: a smooth variety with corners compactifying
$\mathrm{Conf}_n(X\setminus D)$, with boundary stratification indexed
by \emph{rigid planted-forest types}. A planted-forest type $\rho$
consists of a partition of $\{1,\ldots,n\}$ into blocks (the roots),
a rooted tree per block recording the hierarchy of collisions, and a
grid-depth assignment $w_v\in\mathbb Z_{\ge 0}$ at each vertex
recording the order of tangency with~$D$. The ordinary FM is the
case $D=\emptyset$.

The \emph{degeneration formula}: for a semistable $W\to B$ with
singular fibre $W_0=\bigcup Y_v$ and a rigid type $\rho$,
\[
\mathrm{FM}_n(W/B)(\rho)\longrightarrow\prod_{v\in V(S_\rho)}\mathrm{FM}_{I_v}(Y_v|D_v)
\]
is a proper birational correspondence factoring the global FM
geometry into a fibre product of local pieces. Tropicalisation:
$\operatorname{Trop}(\mathrm{FM}_n(\mathbb C|D))=G_{\mathrm{pf}}$, the
planted-forest poset with grafting as partial order. The planted
forest poset is the combinatorial skeleton of the geometric input,
just as the stable graph complex is the skeleton of the Feynman
transform.

The \emph{Mok codimension formula}: for a type with grid depths
$(w_1,\ldots,w_r)$ and tree vertex counts $(|V(T_j)|)_{j=1,\ldots,s}$,
\[
\operatorname{codim}(W_\rho)=\sum_{i=1}^r w_i+\sum_{j=1}^s(|V(T_j)|-1).
\]
For $\mathrm{FM}_3$: four codimension-one strata ($\{12\}$, $\{23\}$,
$\{13\}$, $\{123\}$), three codimension-two strata
($\{12{<}123\}$, $\{23{<}123\}$, $\{13{<}123\}$).

The chain complexes $\cC^{\log\mathrm{FM}}_{X|D}(n):=C_\bullet(\mathrm{FM}_n(X|D))$
carry cooperadic structure: the inclusion of boundary strata defines
cocomposition maps. For a rigid stable graph~$\Gamma$,
\[
\Delta^{\log}_\Gamma:=(\nu_\Gamma)_*\circ\operatorname{Res}_{D^{\log}_\Gamma}\colon\cC^{\log\mathrm{FM}}_{\mathrm{mod}}(g,n)\to\bigotimes_{v\in V(\Gamma)}\cC^{\log\mathrm{FM}}_{\mathrm{mod}}\bigl(g(v),I_v\sqcup H(v)\bigr)[-|E_{\mathrm{int}}(\Gamma)|],
\]
where $\nu_\Gamma$ is the Mok birational correspondence and
$\operatorname{Res}_{D^{\log}_\Gamma}$ is the residue along the
logarithmic boundary divisor.

\subsection*{3.6.\enspace The convolution $L_\infty$-algebra}

A twisting morphism $\alpha\colon\cC\to\cP$ (satisfying
$\partial\alpha+\alpha\star\alpha=0$) determines a convolution
$sL_\infty$-algebra $\operatorname{hom}_\alpha(\cC,\cP)$
(Robert-Nicoud--Wierstra, after Vallette) with underlying space
$\prod_{n\ge 1}\Hom_{\Sigma_n}(\cC(n),\cP(n))$. The brackets
\[
\ell_k(f_1,\ldots,f_k)(x)=\sum\gamma_\cP\bigl((\alpha\circ 1)(f_1\otimes\cdots\otimes f_k)\Delta^{(k)}_\cC(x)\bigr)
\]
come from iterated cooperadic decomposition. $\ell_1$ is the
convolution differential; $\ell_2$ is the Lie bracket of
coderivations; higher $\ell_k$, $k\ge 3$, arise whenever $\cC$ is not
concentrated in a single degree.

The $sL_\infty$-algebra extends to $\infty$-morphisms in either slot
separately but not both simultaneously: the convolution is
functorial in each slot, but does not assemble into a bifunctor. The
MC3 categorical lift proceeds one slot at a time. The bar-cobar
adjunction is a Quillen equivalence (Vallette); different
$\mathrm{Ch}_\infty$ presentations yield quasi-isomorphic
deformation theories.

\subsection*{3.7.\enspace Assembly}

The modular convolution $L_\infty$-algebra:
\begin{equation}\label{eq:pref-convolution}
\mathfrak g^{\mathrm{mod}}_\cA:=\Convinf\!\Bigl(\cC^{\log\mathrm{FM}}_{\mathrm{mod}},\operatorname{End}_{\mathrm{Ch}_\infty}(A^{\mathrm{ch}}_\infty)\Bigr).
\end{equation}
The cooperadic input is the log-FM chain cooperad
$\cC^{\log\mathrm{FM}}_{\mathrm{mod}}$; the operadic input is the
$\mathrm{Ch}_\infty$-endomorphism operad of
$A^{\mathrm{ch}}_\infty$. The $E_1$ cousin
$\mathfrak g^{\mathrm{mod},E_1}_\cA$ uses the ribbon-graph-indexed
cooperad and the ordered bar endomorphism operad.

One-slot convention: fix the log-FM (resp.\ ribbon) cooperad as the
``geometric'' slot; vary the chiral algebra input as the ``algebraic''
slot by $\infty$-morphisms.

The \emph{strict chart}: the coderivation dg Lie algebra
\[
\mathfrak g^{\mathrm{mod},\log}_\cA(\mathfrak U):=\Convstr\!\bigl(\cC^{\log\mathrm{FM}}_{\mathrm{mod}},\operatorname{End}_{\mathrm{Ch}_\infty}(A^{\mathrm{ch}}_\infty)\bigr)
\]
has classical commutator bracket, no higher $L_\infty$ brackets, but
the same Maurer--Cartan moduli. The strict chart suffices for
constructing $\Theta_\cA$ and computing shadows; the full $L_\infty$
structure is needed for transfer, formality, and gauge equivalence.

\subsection*{3.8.\enspace The modular bar functor}

In log-FM coordinates,
\[
\mathrm B^{\log}_{\mathrm{mod}}(\cA)(g,n)=\bigoplus_{\Gamma\in\mathsf{Gr}^{\mathrm{st}}_{g,n}}\left(\bigotimes_{v\in V(\Gamma)}\bar B^{\mathrm{ch}}(A^{\mathrm{ch}}_\infty;\operatorname{In}(v))\otimes C_\bullet(\mathrm{FM}^{\log}_\Gamma)\otimes\orline{E_{\mathrm{int}}(\Gamma)}\right)_{\!\operatorname{Aut}(\Gamma)}.
\]
The ingredients: at each vertex the transferred
$\mathrm{Ch}_\infty$-bar contribution on its incoming half-edges;
at each internal edge the propagator $P_\cA=H_\cA^{-1}$; the geometric
factor $C_\bullet(\mathrm{FM}^{\log}_\Gamma)$; the orientation line
$\orline{E_{\mathrm{int}}(\Gamma)}=\det(\Bbbk E(\Gamma))\otimes\det(H_1(\Gamma;\Bbbk))^{-1}$;
the automorphism quotient.

\subsection*{3.9.\enspace The six-component differential}

The total differential has six components, whose mutual cancellations
constitute $D^2=0$:
\[
D^{\log}_{\mathrm{mod}}=\underbrace{d_{\check C}}_{\text{\v Cech}}+\underbrace{d_{\mathrm{Ch}_\infty}}_{\text{transferred chiral}}+\underbrace{d_{\mathrm{coll}}}_{\text{collision}}+\underbrace{d_{\mathrm{sew}}}_{\text{separating}}+\underbrace{d_{\mathrm{pf}}}_{\text{planted-forest}}+\underbrace{\hbar\Delta}_{\text{genus-raising}}.
\]
$d_{\check C}$ is the \v Cech differential of the cover, acting on
\v Cech indices. $d_{\mathrm{Ch}_\infty}$ is the transferred chiral
differential on each vertex, from secondary Borcherds operations.
$d_{\mathrm{coll}}$ contracts an internal edge by extracting the OPE
residue at the corresponding codimension-one collision stratum of
$\overline{\mathrm{FM}}_n(X)$. The three positive-genus boundary
operators act on a pure tensor $a\otimes\eta$ indexed by $\Gamma$:
\begin{align*}
d_{\mathrm{sew}}(a\otimes\eta)&=\sum_{e\in E_{\mathrm{sep}}(\Gamma)}\pm\mu_e(a)\otimes(\xi_e)_*\operatorname{Res}_{D_e^{\log}}(\eta),\\
d_{\mathrm{pf}}(a\otimes\eta)&=\sum_{F\in\mathsf{PF}^{\mathrm{rig}}(\Gamma)}\pm\mu_F(a)\otimes(\xi_F)_*\operatorname{Res}_{D_F^{\log}}(\eta),\\
\hbar\Delta(a\otimes\eta)&=\hbar\sum_{e\in E_{\mathrm{ns}}(\Gamma)}\pm\operatorname{Tr}_e(a)\otimes(\xi_e)_*\operatorname{Res}_{D_e^{\log}}(\eta).
\end{align*}
Separating edges give the clutching morphism (gluing two lower-genus
curves); non-separating edges give the genus-raising self-gluing
(BV operator); planted-forest strata give the genuinely logarithmic
corrections from Mok. Explicitly,
\[
d_{\mathrm{pf}}=\sum_{\rho\in\mathsf{PF}^{\mathrm{rig}}}\epsilon_\rho(\kappa_\rho)_*\mathrm{pr}_\rho^*\otimes\mu_\rho:
\]
Mok's degeneration formula internalised as a summand of the total
differential.

\subsection*{3.10.\enspace $D^2=0$}

The nilpotence $(D^{\log}_{\mathrm{mod}})^2=0$ holds by three
cancellation mechanisms. \emph{Vertex labels}: $d_{\check C}^{\,2}=0$,
$d_{\mathrm{Ch}_\infty}^{\,2}=0$, and
$\{d_{\check C},d_{\mathrm{Ch}_\infty}\}=0$ by the
$\mathrm{Ch}_\infty$-algebra axioms on $A^{\mathrm{ch}}_\infty$.
\emph{One-edge expansions}: anticommutators of a vertex differential
with an edge-contraction differential vanish by the Leibniz rule for
coderivations. \emph{Codimension-two cancellation}: the sum of all
products of two distinct edge-contraction operators vanishes, because
the codimension-two boundary of Mok's log-FM compactification is a
normal-crossings divisor (Mok, Theorem~3.3.1) and each codimension-two
face appears as the intersection of exactly two codimension-one faces
with opposite orientations.

The third mechanism is the nontrivial one. The normal-crossings
property is essential: without it, codimension-two faces could have
excess intersection, and the cancellation would fail.


% ====================================================================
\section*{4.\quad The universal Maurer--Cartan element}
% ====================================================================

Master conjectures MC1 through MC4 are \emph{proved} (PBW
concentration, MC element existence, thick generation in every type,
completion tower). MC5 is partially proved: the analytic HS-sewing
package is established at all genera, while the genuswise BV/BRST/bar
identification remains conjectural; at genus~$0$ the algebraic
BRST/bar comparison is proved
(Theorem~\ref{thm:algebraic-string-dictionary}) and the tree-level
amplitude pairing is conditional on
Corollary~\ref{cor:string-amplitude-genus0}. Theorem A (bar-cobar
adjunction), Theorem B (bar-cobar inversion on the Koszul locus),
Theorem C (complementarity: C1 unconditional, C2
\textsc{(uniform-weight)}), Theorem~D
(obstruction $\operatorname{obs}_g=\kappa\lambda_g$
\textsc{(uniform-weight)}; multi-weight correction
$\delta F_g^{\mathrm{cross}}$), and Theorem~H (chiral
Hochschild concentration $\{0,1,2\}$, $\dim\le 4$) are the five
$\Sigma_n$-coinvariant projections of a single $E_1$ object.

\subsection*{4.1.\enspace $\Gamma$-amplitudes and Taylor coefficients}

Each stable graph produces a multilinear operation on the convolution
algebra, built from chiral operations at vertices, propagators along
edges, and geometric factors from log-FM strata. For
$\Gamma\in\mathsf{Gr}^{\mathrm{st}}$ with $|V(\Gamma)|=k$ vertices and
elements $f_1,\ldots,f_k\in\gAmod$, the \emph{$\Gamma$-amplitude}
assembles
\[
\ell_\Gamma(f_1,\ldots,f_k)=\mu_\Gamma\circ(\alpha^{\mathrm{mod}}_\Gamma\otimes f_1\otimes\cdots\otimes f_k)\circ\Delta^{\log}_\Gamma,
\]
a $k$-linear map of cohomological degree $2-2g(\Gamma)-k$. Here
$\mu_\Gamma$ is the graph composition (transferred chiral operations
at vertices, propagators along edges), $\alpha^{\mathrm{mod}}_\Gamma$
is the modular twisting morphism component, and $\Delta^{\log}_\Gamma$
is the cooperadic cocomposition.

The Taylor coefficients of the modular $L_\infty$ structure sum over
stable graphs of fixed genus and vertex count:
\[
\ell_k^{(g)}=\sum_{\substack{\Gamma\in\mathsf{Gr}^{\mathrm{st}},\;|V(\Gamma)|=k\\b_1(\Gamma)+\sum_v g(v)=g}}\frac{1}{|\operatorname{Aut}(\Gamma)|}\ell_\Gamma.
\]
The homotopy Jacobi identity is encoded by $D^2=0$ on the total
convolution algebra.

\medskip
\noindent\textbf{Low-order Taylor coefficients.}\enspace
At tree level, the binary bracket is the transferred Lie bracket:
$\ell_2^{(0)}(\alpha,\beta)=\pi[m_2(\iota(\alpha),\iota(\beta))]$.
The ternary bracket sums over the three binary trees in
$\mathsf{Trees}_3$ (channels $s$, $t$, $u$), each inserting a homotopy
$h$ at the internal edge. The quaternary bracket sums over the
fifteen labeled trees in $\mathsf{Trees}_4$ (two unlabeled topologies:
three balanced $((ab)(cd))$ and twelve caterpillar $(((ab)c)d)$),
each with two propagator insertions; this is the classical homological
perturbation lemma in its $L_\infty$ incarnation.

The first genus-raised bracket is unary, the BV operator on the single
graph with one vertex and one self-loop:
\[
\ell_1^{(1)}(\alpha)=\operatorname{tr}(P_\cA\circ\alpha\circ\delta^{\mathrm{ns}})=\sum_i\langle e^i,\alpha(e_i)\rangle_{\mathrm{Sh}},
\]
trace over the propagator contracting two half-edges at the
non-separating node of $\overline{\cM}_{1,1}$. The genus-one binary
bracket
\[
\ell_2^{(1)}(\alpha,\beta)=\operatorname{tr}\bigl(P_\cA\circ m_2(\iota\alpha,-)\circ\delta^{\mathrm{ns}}\circ\iota\beta\bigr)+(\alpha\leftrightarrow\beta)
\]
accounts for the separating node of $\overline{\cM}_{1,2}$. These
low-order brackets are the only data needed for shadow computations
through degree~$4$.

\subsection*{4.2.\enspace The bar-intrinsic construction}

$\Theta_\cA$ requires no equation to solve: it is extracted from a
differential that already squares to zero.

Decompose the total bar differential
$D_\cA\colon\mathrm B^{\mathrm{mod}}(\cA)\to\mathrm B^{\mathrm{mod}}(\cA)$
as $D_\cA=d_0+\Theta_\cA$, where
$d_0=d_{\mathrm{int}}+[\tau,-]$ is the tree-level local part. The
\emph{universal Maurer--Cartan element} is the remainder,
\begin{equation}\label{eq:pref-theta}
\Theta_\cA:=D_\cA-d_0=\sum_{\Gamma\in\mathsf{Gr}^{\mathrm{st}}}\frac{W^{\log\mathrm{FM}}_\Gamma}{|\operatorname{Aut}(\Gamma)|}\otimes\Phi^\cA_\Gamma\in\operatorname{MC}(\gAmod),
\end{equation}
with $W^{\log\mathrm{FM}}_\Gamma\in H^*(\overline{\cM}^{\log}_{g(\Gamma),n(\Gamma)})$
the log-FM fundamental class of the stratum and
$\Phi^\cA_\Gamma\in\Hom(\bar\cA^{\otimes n(\Gamma)},\Bbbk)$ the chiral
amplitude (propagators on internal edges, operations at vertices,
cyclic trace at the output).

The MC property is immediate. $D_\cA^2=0$ by the normal-crossings
property of $\overline{\cM}_{g,n}^{\log}$; $d_0^2=0$ by construction;
expanding $(d_0+\Theta_\cA)^2=0$ gives
\[
[d_0,\Theta_\cA]+\tfrac12[\Theta_\cA,\Theta_\cA]=0.
\]
The MC property is a formal consequence of $\partial^2=0$ on the
moduli space of curves, requiring no simple-Lie restriction, no
convergence hypothesis, and no choice of basis. This is MC2
(conjecture-turned-theorem) in its cleanest form.

The $E_1$ analogue $\Theta^{E_1}_\cA$ is constructed identically on
the ordered, ribbon-graph-indexed convolution algebra, and
$\Theta_\cA=\mathrm{av}(\Theta^{E_1}_\cA)$ under the coinvariant map
$F\!\Ass\to F\Com$.

\subsection*{4.3.\enspace $\Theta_\cA$ as universal modular twisting morphism}

The Maurer--Cartan space of the modular convolution $L_\infty$
algebra is identified with the space of modular twisting morphisms,
\[
\operatorname{MC}_\bullet(\gAmod)\simeq\operatorname{Tw}_{\alpha^{\mathrm{mod}}}\Bigl(\cC^{\log\mathrm{FM}}_{\mathrm{mod}},\operatorname{End}_{\mathrm{Ch}_\infty}(A^{\mathrm{ch}}_\infty)\Bigr).
\]
Under this identification $\Theta_\cA$ is the \emph{universal modular
twisting morphism}, the operadic analogue of the canonical flat
connection on a principal bundle. A twisting morphism
$\alpha\colon\cC\to\cP$ determines a twisted composition product
$\cC\circ_\alpha\cP$ whose differential squares to zero iff $\alpha$
satisfies the MC equation. $\Theta_\cA$ is the unique twisting
morphism that simultaneously twists the modular bar complex and
intertwines with the log-FM boundary structure.

The principal-bundle analogy is exact. The modular bar bundle
$\mathrm B^{\log}_{\mathrm{mod}}(\cA)\to\overline{\cM}_{g,n}$ is the
principal bundle; $\Theta_\cA$ is the connection; the MC equation is
flatness; the shadow tower (\S6) is the Chern--Weil theory. The
structure group is the pro-unipotent MC gauge group
$\exp(\mathfrak g^{\mathrm{mod},0}_\cA)$.

\subsection*{4.4.\enspace Tridegree, depth, and convergence}

Each coefficient of $\Theta_\cA$ carries a natural tridegree
\[
(g,n,d)=(\text{loop genus},\text{external degree},\text{log boundary depth}),
\]
with $d(\rho)=\#E_{\mathrm{sep}}(\Gamma_\rho)+\#E_{\mathrm{ns}}(\Gamma_\rho)+\operatorname{depth}_{\mathrm{pf}}(\Gamma_\rho)$
the codimension of the log-FM stratum. The depth filtration
$F^d\gAmod=\bigoplus_{d'\ge d}(\gAmod)_{*,*,d'}$ makes the convolution
algebra complete:
\[
\gAmod=\varprojlim_d\gAmod/F^d\gAmod.
\]
The inverse limit $\Theta_\cA=\varprojlim_r\Theta_\cA^{\le r}$
converges unconditionally: no analytic input, no growth estimate, no
restriction on the algebra. Convergence is purely algebraic,
guaranteed by completeness of the depth filtration. This is MC4
(completion tower): proved.

At fixed $(g,n)$, only finitely many stable graph types contribute.
The depth filtration refines this further: at depth $d$, the stratum
has codimension $d$ in $\overline{\cM}_{g,n}^{\log}$, and the
amplitude has internal-edge count bounded by $d$.

\subsection*{4.5.\enspace The weight filtration}

Decompose the MC equation by internal-edge count,
\[
\Theta_\cA=\sum_{r\ge 1}\Theta^{[r]}_\cA,\qquad\Theta^{[r]}_\cA=\sum_{\substack{\Gamma\in\mathsf{Gr}^{\mathrm{st}}\\|E_{\mathrm{int}}(\Gamma)|=r}}\frac{W^{\log\mathrm{FM}}_\Gamma}{|\operatorname{Aut}(\Gamma)|}\otimes\Phi^\cA_\Gamma.
\]
Weight~$1$: $[D_{\mathrm{loc}},\Theta^{[1]}_\cA]=0$. Weight~$N$:
$[D_{\mathrm{loc}},\Theta^{[N]}_\cA]+\tfrac12\sum_{i+j=N}[\Theta^{[i]}_\cA,\Theta^{[j]}_\cA]=0$.
The obstruction to extending from weight $N-1$ to weight $N$ lies in
$H^2(\gAmod/F^{N+1},D_{\mathrm{loc}})$. The bar-intrinsic
construction solves all stages simultaneously.

\medskip
\noindent\textbf{Weight~$1$: boundary divisors.}\enspace
In degree $(0,4)$: the three boundary points of
$\overline{\cM}_{0,4}\cong\mathbb P^1$ correspond to the three
channels $(12|34)$, $(13|24)$, $(14|23)$,
\[
\Theta^{[1]}_\cA\big|_{(0,4)}=[\delta_s]\otimes\Phi_s+[\delta_t]\otimes\Phi_t+[\delta_u]\otimes\Phi_u.
\]
The weight-one equation $\Phi_s+\Phi_t+\Phi_u=F_3$ identifies the
Jacobiator homotopy with a boundary in the bar complex: the first
Borcherds identity. In degree $(1,1)$: the single boundary divisor
$\delta_{\mathrm{irr}}$ gives
$\Theta^{[1]}_\cA|_{(1,1)}=[\delta_{\mathrm{irr}}]\otimes\Delta$ with
$\Delta=\sum_i e_i\otimes e^i$ the Casimir. In degree $(0,5)$: the
ten boundary divisors of $\overline{\cM}_{0,5}$ (Petersen graph)
recover the pentagon identity. The entire $A_\infty$ hierarchy is
the tree-level projection of a single MC element.

\medskip
\noindent\textbf{Weight~$2$: the first planted-forest layer.}\enspace
At weight~$2$, codimension-two strata enter: classical products of
two boundary divisors plus genuinely logarithmic planted-forest
strata with no DM counterpart. The weight-two equation is
$[D_{\mathrm{loc}},\Theta^{[2]}_\cA]=-\tfrac12[\Theta^{[1]},\Theta^{[1]}]$,
and the planted-forest correction $\Theta^{[2]}_{\mathrm{pf}}$ records
how collision limits in $\mathrm{FM}_n(X|D)$ correct the naive square.
For the $\mathcal W_3$ central-parameter line, cubic gauge triviality
kills the weight-$2$ obstruction; the first genuine invariant appears
at weight~$3$.

\subsection*{4.6.\enspace The all-genus recursion}

Projecting the MC equation to fixed $(g,n)$,
\[
d\Theta^{(g,n)}_\cA+\tfrac12\!\!\!\sum_{\substack{g_1+g_2=g\\n_1+n_2=n+2}}\!\!\!\ell_2^{(0)}(\Theta^{(g_1,n_1)}_\cA,\Theta^{(g_2,n_2)}_\cA)+\Delta_{\mathrm{ns}}\Theta^{(g-1,n+2)}_\cA+\sum_{k\ge 3}\frac{1}{k!}\ell_k^{(\mathrm{tree})}(\Theta_\cA,\ldots,\Theta_\cA)^{(g,n)}=0.
\]
Four terms: the internal differential, separating clutching,
non-separating BV, and planted-forest corrections. This recursion
determines $\Theta^{(g,n)}_\cA$ inductively from lower $(g',n')$ in
lexicographic order. The base cases are
$\Theta^{(0,3)}_\cA$ (the chiral bracket) and $\Theta^{(0,4)}_\cA$
(the Jacobiator data). The recursion terminates at each $(g,n)$.

\medskip
\noindent\textbf{The clutching identity.}\enspace
On each rigid boundary stratum~$\rho$,
\[
\Theta_\cA\big|_\rho=(\kappa_\rho)_*\!\left(\bigotimes_{v\in V(S_\rho)}\Theta_{\cA,v}\right).
\]
The amplitude on the degenerate curve equals the product of amplitudes
on the irreducible components, sewn by propagators along the nodes.
This is the MC-level sewing axiom, and the starting point for the
clutching law of the quartic shadow (\S6).

\bigskip
\noindent
Every invariant in this monograph is a projection of $\Theta_\cA$:
the curvature $\kappa$, the spectral discriminant $\Delta_\cA$, the
shadow metric $Q_L$, the sewing amplitudes, the $\hat A$-genus, the
Fredholm determinant, and the quartic contact invariant
$Q^{\mathrm{contact}}$. The five main theorems are five projections.

\medskip
\begin{center}
\small
\renewcommand{\arraystretch}{1.15}
\begin{tabular}{lll}
\textbf{Projection of $\Theta_\cA$} & \textbf{Invariant} & \textbf{Section}\\
\hline
Scalar ($r=2$, genus tower) & $\kappa(\cA)$, $F_g$, $\hat A$-genus & \S2\\[1pt]
Binary genus-$0$ residue & classical $r$-matrix, Yangian seed & \S7.7\\[1pt]
Binary Verdier dual & complementarity $Q_g\oplus Q_g^!$ & \S2\\[1pt]
Bar-cobar counit & inversion $\Omega(\barB(\cA))\xrightarrow{\sim}\cA$ & \S1\\[1pt]
Chiral Hochschild & $\operatorname{ChirHoch}^*(\cA)$, deformation ring & \S9\\[1pt]
Spectral discriminant & $\Delta_\cA(x)$, genus-$1$ Hessian & \S5\\[1pt]
Cubic shadow ($r=3$) & gauge-trivial obstruction & \S6\\[1pt]
Quartic resonance ($r=4$) & $\mathfrak Q^{\mathrm{contact}}$, clutching law & \S6\\[1pt]
MC tautological descent & shadow CohFT, $R^*(\overline{\cM}_{g,n})$ & \S5\\[1pt]
Primitive kernel & shell equations, branch BV & \S5\\[1pt]
Dirichlet--sewing lift & $S_\cA(u)$, Euler--Koszul class & \S8\\[1pt]
Swiss-cheese colour decomposition & $\alpha_T$, six projections (Vol.~II) & \S10
\end{tabular}
\end{center}
\medskip

\noindent
These five theorems are five projections of $\Theta_\cA$. Theorem~A
constructs the arena in which $\Theta_\cA$ lives. Theorem~B
recovers $\cA$ from $\Theta_\cA$ on the Koszul locus. Theorem~C
splits the genus tower into Lagrangian summands (C1 unconditional, C2
\textsc{(uniform-weight)}). Theorem~D extracts the scalar $\kappa$
with $\hat A$-genus generating function (uniform-weight; multi-weight
adds $\delta F_g^{\mathrm{cross}}$). Theorem~H identifies the
deformation ring via $\operatorname{ChirHoch}^*$ concentration.

The finite-order projections form an obstruction tower
$\Theta_\cA^{\le 2}\to\Theta_\cA^{\le 3}\to\Theta_\cA^{\le 4}\to\cdots$
whose generating function is algebraic of degree~$2$. On each
primary line,
$H(t)=t^2\sqrt{Q_L(t)}$ is determined by three seed invariants
$(\kappa,\alpha,S_4)$. The inverse limit exists and is automatically
Maurer--Cartan because $D_\cA^2=0$ (MC4).

In three dimensions, $\Theta_\cA$ extends to an open/closed
Maurer--Cartan element
$\Theta^{\mathrm{oc}}_\cA=\Theta_\cA+\sum_j\mu^{M_j}$
packaging the holomorphic-topological QFT on
$\mathbb C_z\times\mathbb R_t$: the bar complex $B(\cA)$, coassociative over $(\mathrm{ChirAss})^!$,
supplies the holomorphic factorization data; the derived center pair
$(\cA, \cZ^{\mathrm{der}}_{\mathrm{ch}}(\cA))$ carries the
$\SCchtop$-algebra structure encoding both closed and open colours;
and $\cZ^{\mathrm{der}}_{\mathrm{ch}}(\cA)$ is the universal bulk
(Volume~II).

The bar complex thus organises five mathematical structures as
projections of one object: the genus expansion
$F_g=\kappa\cdot\lambda_g^{\mathrm{FP}}$ \textsc{(uniform-weight)} is
a GUE matrix model at $N^2=\kappa$; the ordered bar
$\barB^{\mathrm{ord}}(\cA)$ produces $U_q(\mathfrak g)$ via FRT
reconstruction; the quantum-group $R$-matrix yields knot invariants;
the shadow metric $Q_L$ defines a Bridgeland stability condition with
universal wall phase $\pi/4$; and the Stokes constants
$S_n/S_1=(-1)^{n+1}n$ encode the resurgent structure.

\medskip
\noindent\textbf{Corollary} (MC replaces the sum over topologies).
On the Koszul locus, the tree-level complex carries a contracting
homotopy (from the PBW degeneration), so $\Theta^{(g)}$ is determined
up to gauge-equivalent MC elements by induction from $\Theta^{(0)}$.
No free parameters appear at any genus; distinct choices of
contracting homotopy yield gauge-equivalent MC elements; no choice of
bulk topologies is required. The genus expansion is computed
algebraically from the genus-$0$ chiral bracket by the clutching
recursion of \S4.6, not approximated perturbatively. The
Maloney--Witten problem (negative Fourier coefficients in the
Poincar\'e series) does not arise: the MC equation determines each
$F_g$ uniquely by algebraic induction, without summing over bulk
topologies at fixed genus. The non-perturbative sum over distinct
three-manifold topologies lies outside the MC framework.


% ====================================================================
\section*{5.\quad The modular tangent complex and Chern--Weil theory}
% ====================================================================

The MC element $\Theta_\cA$ exists unconditionally and satisfies the
MC equation. Its deformation theory is controlled by the modular
tangent complex, obtained by twisting the convolution algebra by
$\Theta_\cA$ itself.

\subsection*{5.1.\enspace The $L_\infty$-twisted differential}

The MC element twists the convolution differential. For $x\in\gAmod$,
\[
d_{\Theta_\cA}(x):=\sum_{n\ge 0}\sum_{g\ge 0}\frac{\hbar^g}{n!}\ell_{n+1}^{(g)}(\Theta_\cA^{\otimes n},x)=d(x)+[\Theta_\cA,x]+\tfrac{1}{2!}\ell_3^{(0)}(\Theta_\cA,\Theta_\cA,x)+\hbar\ell_1^{(1)}(x)+\cdots
\]
In the strict chart this collapses to $d_{\Theta_\cA}(x)=d(x)+[\Theta_\cA,x]$.
The identity $d_{\Theta_\cA}^2=0$ is the MC equation for $\Theta_\cA$.

\subsection*{5.2.\enspace The modular tangent complex}

The \emph{modular tangent complex at $\Theta_\cA$}:
\[
T^{\mathrm{mod}}_{\Theta_\cA}:=(\gAmod,d_{\Theta_\cA}).
\]
Its cohomology governs deformations: $H^0$ infinitesimal automorphisms
(stabiliser), $H^1$ first-order deformations (tangent space), $H^2$
obstructions. $H^2=0$ makes the MC moduli smooth at $\Theta_\cA$;
$H^1=0$ makes $\Theta_\cA$ rigid.

For a one-parameter family $\cA_t$ varying in $k(t)$ or $c(t)$, the
tangent complex forms a local system over the parameter space and
the obstruction class traces a path in $H^2$. The vanishing locus
is the moduli of liftable deformations.

\subsection*{5.3.\enspace Three characteristic projections}

Three successive projections of $\Theta_\cA$ capture progressively
finer structure: $\kappa(\cA)$ (leading scalar curvature),
$\Delta_\cA(x)$ (characteristic polynomial of a finite-rank operator),
$\mathfrak R^{\mathrm{mod}}_4(\cA)$ (first nonlinear characteristic
class).

\medskip
\noindent\textbf{The scalar curvature.}\enspace
\[
\kappa(\cA)=\operatorname{tr}(\Theta_\cA)_{2,0}=\operatorname{tr}(\pi\circ m_2\circ\iota^{\otimes 2}\circ\Delta_{\mathrm{sep}})\in\Bbbk.
\]
This is the first Chern class of the modular bar bundle:
$\operatorname{obs}_g(\cA)=\kappa(\cA)\cdot\lambda_g$
\textsc{(uniform-weight)}; multi-weight at $g\ge 2$ adds
$\delta F_g^{\mathrm{cross}}$. The Mumford relation
$\lambda_g^{g+1}=0$ confines genus-$g$ scalar obstructions to
dimension at most~$g$.

Explicit values across the standard landscape:
\begin{center}
\renewcommand{\arraystretch}{1.2}
\begin{tabular}{lccl}
\textbf{Algebra $\cA$} & $\kappa(\cA)$ & $\kappa(\cA^!)$ & \textbf{Sum rule}\\
\hline
Heisenberg $\cH_k$ & $k$ & $-k$ & $0$\\[1pt]
Affine $\widehat{\fg}_k$ & $\tfrac{(k+h^\vee)\dim\fg}{2h^\vee}$ & $-\tfrac{(k+h^\vee)\dim\fg}{2h^\vee}$ & $0$\\[4pt]
Virasoro $\mathrm{Vir}_c$ & $c/2$ & $(26-c)/2$ & $13$\\[1pt]
$\beta\gamma$ system & $1$ & $-1$ & $0$\\[1pt]
Lattice $V_\Lambda$ & $\operatorname{rank}(\Lambda)$ & $-\operatorname{rank}(\Lambda)$ & $0$\\[1pt]
$\mathcal W_N(\widehat{\mathfrak{sl}}_N)$ & $c\cdot(H_N-1)$ & $c'\cdot(H_N-1)$ & $K_N\cdot(H_N-1)$\\[4pt]
Free fermion $\psi$ & $1/4$ & $-1/4$ & $0$
\end{tabular}
\end{center}
$H_N=\sum_{j=1}^N 1/j$ is the harmonic number. The sum rule
$\kappa(\cA)+\kappa(\cA^!)=0$ holds for Kac--Moody, free fields,
lattice algebras. For $\mathcal W$-algebras the sum is a nonzero
constant: $\kappa(\mathrm{Vir}_c)+\kappa(\mathrm{Vir}_{26-c})=13$ and
$\kappa(\mathcal W_{N,c})+\kappa(\mathcal W_{N,c'})=K_N(H_N-1)$ with
$K_N=c+c'$. Self-duality occurs at $c=13$ for Virasoro and $c=50$
for $\mathcal W_3$.

\medskip
\noindent\textbf{The spectral discriminant.}\enspace
The Koszul dual Hilbert series
\[
h_{\cA^!}(q)=\sum_{n\ge 0}\dim H^n(\barB(\cA))\,q^n
\]
is rational for every standard family: it satisfies a holonomic
recurrence of finite order $d$, and the transfer matrix
$T_{\mathrm{rec}}$ is a $d\times d$ matrix. Define the
\emph{spectral branch object} as a finite-rank perfect complex
$V^{\mathrm{br}}_\cA$ with \emph{branch transport operator}
$T_{\mathrm{br},\cA}\colon V^{\mathrm{br}}_\cA\to V^{\mathrm{br}}_\cA$
and
\begin{equation}\label{eq:pref-delta-def}
\Delta_\cA(x):=\operatorname{sdet}(1-xT_{\mathrm{br},\cA}).
\end{equation}
The rank of $V^{\mathrm{br}}_\cA$ equals the degree of algebraicity of
$h_{\cA^!}$. Three properties: $\Delta_\cA$ is self-dual
($\Delta_{\cA^!}=\Delta_\cA$), unlike $\kappa$ which is anti-symmetric
for KM/free; exact factorization functors preserving quadratic OPE
data preserve $T_{\mathrm{br}}$; the spectral radius equals
$\dim\mathfrak g$ for affine KM at generic level.

For $\mathfrak{sl}_3$: the degree-$3$ recurrence has transfer matrix
eigenvalues $8$, $(3\pm\sqrt{13})/2$, and
$\Delta_{\widehat{\mathfrak{sl}}_3}(x)=(1-8x)(1-3x-x^2)$. For
Heisenberg: $h_{\mathcal H_k^!}(q)=1+q$ and $\Delta_{\mathcal H_k}(x)=1$.
For $\beta\gamma$: $h=\sqrt{(1+q)/(1-3q)}$ and
$\Delta_{\beta\gamma}(x)=(1-3x)(1+x)$.

In the Chern--Weil dictionary, $\Delta_\cA$ plays the role of the
Chern character:
$\operatorname{ch}_r(\cA)=\tfrac{1}{r!}\tfrac{d^r}{dt^r}|_{t=0}[-\log\Delta_\cA(t)]=\tfrac{1}{r}\operatorname{tr}(T_{\mathrm{br}}^r)$.

\medskip
\noindent\textbf{The quartic shadow.}\enspace
\[
\mathfrak R^{\mathrm{mod}}_{4,g,n}(\cA)=\operatorname{pr}_{4,g,n}(\Theta_\cA)\in H^*(\overline{\cM}_{g,n})\otimes(\bar\cA^\vee)^{\otimes n}_{\mathrm{cyc}}.
\]
This is the first \emph{nonlinear} characteristic class: the first
invariant not recoverable from $\kappa$ and $\Delta_\cA$. Its
genus-zero, four-point specialisation is the quartic contact invariant
\[
Q^{\mathrm{contact}}_{\mathrm{Vir}}=\frac{10}{c(5c+22)},\qquad Q^{\mathrm{contact}}_{\cH}=0,\qquad Q^{\mathrm{contact}}_{\beta\gamma}\neq 0.
\] For Heisenberg the
quartic vanishes (no cubic OPE pole). For $\beta\gamma$ the quartic
is the first gauge-nontrivial shadow (the cubic is gauge-removable).
For Virasoro, both cubic and quartic are nonvanishing.

\subsection*{5.4.\enspace The Chern--Weil dictionary}

The analogy between $\Theta_\cA$ and a connection $1$-form is exact.
The modular bar bundle
$\mathrm B^{\log}_{\mathrm{mod}}(\cA)\to\overline{\cM}_{g,n}$ is a
principal bundle for the gauge group of the convolution algebra;
$\Theta_\cA$ is its flat connection (flatness automatic); the shadow
tower is its Chern--Weil theory. Holonomy becomes factorisation
homology, gauge equivalence becomes MC gauge, and the Weil
homomorphism becomes the shadow algebra
$\cA^{\mathrm{sh}}=H_\bullet(\Defcyc^{\mathrm{mod}})$.

\begin{center}
\renewcommand{\arraystretch}{1.25}
\begin{tabular}{ll}
\textbf{Classical Chern--Weil} & \textbf{Modular factorization}\\
\hline
Principal $G$-bundle $P\to M$ & Modular bar bundle $\mathrm B^{\log}_{\mathrm{mod}}(\cA)\to\overline{\cM}_{g,n}$\\[1pt]
Connection $\omega\in\Omega^1(P,\fg)$ & MC element $\Theta_\cA\in\operatorname{MC}(\gAmod)$\\[1pt]
Curvature $\Omega=d\omega+\tfrac12[\omega,\omega]$ & $D\Theta_\cA+\tfrac12[\Theta_\cA,\Theta_\cA]=0$\\[1pt]
Flatness $\Omega=0$ & Automatic (bar-intrinsic)\\[1pt]
First Chern $c_1$ & $\kappa(\cA)=\operatorname{tr}(\Theta_\cA)_{2,0}$\\[1pt]
Chern character & Spectral determinant $\Delta_\cA(t)$\\[1pt]
Higher Chern classes $c_r$ & Shadows $\mathfrak R^{\mathrm{mod}}_{4,g,n}$, $\operatorname{Sh}_r$ ($r\ge 5$)\\[1pt]
Weil homomorphism & Shadow algebra $\cA^{\mathrm{sh}}=H_\bullet(\Defcyc^{\mathrm{mod}})$\\[1pt]
$\operatorname{ad}$-invariant polynomial & Cyclic cocycle in $\Defcyc^{\mathrm{mod}}(\cA)$\\[1pt]
Cyclic pairing $\operatorname{tr}(XY)$ & Verdier duality on $\operatorname{Ran}(X)$\\[1pt]
Group multiplication & Operadic composition in $\cP^{\mathrm{ch}}$\\[1pt]
Coproduct on $\cO(G)$ & Operadic cocomposition / clutching\\[1pt]
Lagrangian splitting $T^*G=\fg\oplus\fg^*$ & Complementarity $Q_g(\cA)\oplus Q_g(\cA^!)$\\[1pt]
Holonomy $\operatorname{Hol}(\gamma)$ & Factorisation homology $\int_\Sigma\cA$\\[1pt]
Gauge equivalence & MC gauge in $\operatorname{MC}_\bullet(\gAmod)$\\[1pt]
Transgression $\operatorname{CS}(\omega)$ & Shadow obstruction tower (\S6)\\[1pt]
Flat sections $\nabla s=0$ & Twisted cochains $d_{\Theta_\cA}(x)=0$
\end{tabular}
\end{center}

\subsection*{5.5.\enspace The genus spectral sequence}

The genus filtration $F^g L_{\mathrm{mod}}:=\bigoplus_{g'\ge g}(\gAmod)_{g',*,*}$
yields a convergent spectral sequence. Write
$D_0=D_{\mathrm{int}}+D_{\mathrm{sep}}$ for the genus-preserving part
and $D_1=D_{\mathrm{nsep}}$ for the genus-raising part; these satisfy
$D_0^2=0$, $\{D_0,D_1\}=0$, $D_1^2=0$ (a bicomplex).

The spectral sequence: $E_0^{p,q}=\operatorname{gr}_F^p L_{\mathrm{mod}}^{p+q}$
with $d_0=[D_0,-]$; $E_1^{p,*}$ is the associated graded cohomology,
genus by genus; the differential $d_r$ at page $r$ encodes the
obstruction $\operatorname{Ob}_r$; $E_\infty^{p,q}\cong\operatorname{gr}_F^p H^{p+q}(L_{\mathrm{mod}})$.
Convergence follows from completeness of the depth filtration and
pronilpotence of the weight filtration; the $\varprojlim^1$ term
vanishes by Mittag-Leffler.

This spectral sequence is distinct from the PBW spectral sequence
(which decomposes by conformal weight within fixed genus). PBW
computes \emph{what} bar cohomology is; the genus spectral sequence
determines \emph{whether} the MC element extends across genera. For
Heisenberg it degenerates at $E_1$. For Virasoro, $d_1\neq 0$ and
nontrivial differentials persist at every page.

\subsection*{5.6.\enspace Homotopy invariance and functoriality}

The shadow algebra $\cA^{\mathrm{sh}}=H_\bullet(\Defcyc^{\mathrm{mod}}(\cA))$
is a homotopy invariant. For any $\infty_\alpha$-quasi-isomorphism
$f\colon\cA\xrightarrow{\sim}\cA'$,
\[
f_*\colon\operatorname{MC}_\bullet(\gAmod)\xrightarrow{\;\sim\;}\operatorname{MC}_\bullet(\mathfrak g^{\mathrm{mod}}_{\cA'})
\]
is a weak homotopy equivalence. The scalar $\kappa$, the spectral
determinant $\Delta_\cA$, and the full shadow obstruction tower
$(\kappa,\Delta,\mathfrak C,\mathfrak Q,\operatorname{Sh}_5,\ldots)$
consist of quasi-isomorphism invariants; $\Theta_\cA$ itself is a
quasi-isomorphism invariant up to gauge. The assignment
$\cA\mapsto(\gAmod,\Theta_\cA)$ is a functor from chiral algebras with
$\infty_\alpha$-quasi-isos inverted to MC moduli problems. The
convolution $L_\infty$-algebra is functorial in each slot separately
but not as a bifunctor; functoriality of $\cA\mapsto\gAmod$ proceeds
through the source slot.

\subsection*{5.7.\enspace The shadow CohFT and topological recursion}

Projected to $\overline{\cM}_{g,n}$, the MC equation produces
tautological classes and tautological relations. The MC element
determines shadow tautological classes
$\tau_{g,n}(\cA)\in R^*(\overline{\cM}_{g,n+1})$ assembling into a
cohomological field theory (Theorem~\ref{thm:shadow-cohft}): the
factorisation axiom holds by the clutching law, the
$\Sigma_n$-equivariance by the symmetry of stable-graph summation,
the flat identity conditionally on the vacuum lying in~$V$.

Projecting the MC equation to the $(g,n)$-component yields a
tautological relation in $R^*(\overline{\cM}_{g,n})$ at each $(g,n)$:
\[
\sum_{\mathrm{sep}/\mathrm{nsep}}\xi_{\mathrm{sep},*}\bigl(\tau_{g_1,|S|}(\cA)\cdot\tau_{g_2,|S^c|}(\cA)\bigr)+\delta_{\mathrm{pf}}^{(g,n)}=0.
\]

\medskip
\noindent\textbf{WDVV and Mumford from MC.}\enspace
At genus~$0$, degree~$3$: the MC tautological relation on
$\overline{\cM}_{0,4}$ reduces to WDVV. At degree~$2$, genus~$g\ge 2$:
substituting $\tau_{g,2}(\cA)=\kappa\cdot\pi^*\lambda_g$
\textsc{(uniform-weight)} recovers the Mumford $\lambda$-class
clutching relation on $\overline{\cM}_g$.

The complementarity propagator
$P_\cA\colon H^*(\cA)\otimes H^*(\cA^!)\to\mathbb C$ is the Givental
$R$-matrix of the shadow CohFT; when it carries a flat unit (all
standard families with vacuum in $V$), Teleman reconstruction
recovers $\{\tau_{g,n}\}$ at all $(g,n)$ from genus-$0$ data and
$P_\cA$. Eynard--Orantin topological recursion is the MC equation
applied to the shadow obstruction tower.

\begin{center}
\small
\renewcommand{\arraystretch}{1.15}
\begin{tabular}{ll}
\textbf{MC datum} & \textbf{Tautological output}\\
\hline
$\Theta_\cA\in\operatorname{MC}(\gAmod)$ & shadow CohFT $\Omega_{g,n}^\cA$\\[1pt]
$D^2=0$ on $\mathrm{FM}^{\log}$ & CohFT axioms\\[1pt]
$D\Theta+\tfrac12[\Theta,\Theta]=0$ at $(g,n)$ & tautological relation\\[1pt]
$[\Theta,\Theta]|_{\mathrm{sep}}$ & WDVV / associativity\\[1pt]
$\hbar\Delta(\Theta)|_{\mathrm{nsep}}$ & Mumford $\lambda$-class relations\\[1pt]
$P_\cA$ & Givental $R$-matrix\\[1pt]
MC recursion & EO topological recursion
\end{tabular}
\end{center}

\medskip
\noindent\textbf{The Pixton shadow bridge.}\enspace
For class-M algebras, $\Theta_\cA$ generates tautological classes at
every genus. At genus~$2$--$3$, the relations
$\tau_{g,n}(\mathrm{Vir}_c)$ for generic $c$ produce elements of the
Pixton ideal in $R^*(\overline{\cM}_{g,n})$: Pixton's tautological
ring relations arise as shadows of the Virasoro MC equation. On the
semisimple locus, MC-descended relations together with Mumford
relations generate the Pixton ideal
(Theorem~\ref{thm:pixton-from-mc-semisimple}); the non-semisimple
frontier remains open.

\medskip
\noindent\textbf{The B\"ocherer bridge.}\enspace
At genus~$2$, the shadow tautological class $\tau_{2,0}(\mathrm{Vir}_c)$
for lattice vertex algebras computes a central $L$-value via the
B\"ocherer conjecture (Furusawa--Morimoto). For the Leech lattice,
$\chi_{12}=\operatorname{SK}(f_{22})$, and the MC-derived second
Fourier coefficient $c_2\approx-1.918\times 10^{-6}\neq 0$ is the
first nonzero central $L$-value produced by the shadow obstruction
tower.

\subsection*{5.8.\enspace The primitive kernel and cofree reduction}

$\Theta_\cA$ lives in the cofree conilpotent coalgebra underlying the
modular bar construction. By Milnor--Moore for modular operads,
$\Theta_\cA$ is determined by its restriction to primitive corollas
and rigid cutting strata, the \emph{primitive logarithmic modular
kernel}
\[
\mathfrak K_\cA=\sum_{g\ge 0,n\ge 2}K_{g,n}^\cA+\sum_{\rho\in\mathrm{Rig}}K_\rho^\cA,
\]
and the equivalence (Theorem~\ref{thm:primitive-to-global-reconstruction})
\[
(Q_\cA^{\mathrm{mod}})^2=0\iff d\mathfrak K_\cA+\mathfrak K_\cA\star\mathfrak K_\cA+\hbar\Delta_{\mathrm{ns}}(\mathfrak K_\cA)=0,
\]
with $\star$ the primitive pre-Lie product from log-FM residues,
replaces the genus-by-genus obstruction programme by a single equation
on primitive data.

\medskip
\noindent\textbf{Shell equations.}\enspace
Projecting to genus~$2$ and~$3$:
\begin{align*}
R_2(\mathfrak K_\cA)&=\Delta_{\mathrm{ns}}(K_1)+\tfrac12[K_1,K_1]+\sum_{\rho\in\mathrm{Rig}_2}R_\rho^\cA(K_{0,2}^\cA),\\
R_3(\mathfrak K_\cA)&=\Delta_{\mathrm{ns}}(K_2)+[K_1,K_2]+\tfrac{1}{3!}\ell_3(K_1,K_1,K_1)+\sum_\rho R_\rho^\cA(\ldots).
\end{align*}
For Heisenberg, only $\Delta_{\mathrm{ns}}$ survives. For affine
Kac--Moody, $[K_1,K_1]$ contributes the cubic shadow. For
$\beta\gamma$, cubic gauge triviality kills the bracket and the rigid
stratum contributes the quartic contact invariant. For Virasoro and
$\mathcal W_N$, all three terms are active at every genus.

\medskip
\noindent\textbf{The branch BV action.}\enspace
The finite-rank branch-active quotient
$V_\cA^{\mathrm{br}}=V_\cA^+\oplus V_\cA^-$ carries a BV master action
\[
S_\cA^{\mathrm{br}}(x,\xi)=\tfrac12\Omega_\cA(K_\cA^{\mathrm{br}}x,x)+\sum_{m+n\ge 3}\frac{1}{m!n!}\mu_{m,n}^\cA(x^{\otimes m};\xi^{\otimes n}),
\]
satisfying
$\hbar\Delta_\cA^{\mathrm{br}}S_\cA^{\mathrm{br}}+\tfrac12\{S_\cA^{\mathrm{br}},S_\cA^{\mathrm{br}}\}_{\mathrm{BV}}=0$.
The metaplectic half-density $\delta_\cA\in 1+x\Bbbk[[x]]$ satisfies
$\delta_\cA^2=\Delta_\cA=\det(1-xT_\cA^{\mathrm{br,red}})$: the
spectral discriminant is the square of a canonical half-density.

\subsection*{5.9.\enspace Entanglement from the modular characteristic}

The bar coalgebra carries a natural bipartite structure: its coproduct
splits a configuration of points into two halves, one on each side of
a spatial cut. For a spatial bipartition of an interval of length $L$
with UV cutoff $\varepsilon$, the scalar projection of $\Theta_\cA$
restricted to the bipartite cut yields
\[
S_{\mathrm{EE}}(\cA)=\frac{c(\cA)}{3}\log\frac{L}{\varepsilon}.
\]
For Virasoro $c=2\kappa$; the relationship between $c$ and $\kappa$
is family-dependent. Complementarity gives the entanglement sum rule
$S_{\mathrm{EE}}(\cA)+S_{\mathrm{EE}}(\cA^!)=\tfrac{c(\cA)+c(\cA^!)}{3}\log(L/\varepsilon)$;
for the Virasoro family $\mathrm{Vir}_c^!=\mathrm{Vir}_{26-c}$ this
gives $c+c'=26$ and the universal sum $26/3\cdot\log(L/\varepsilon)$.

The four shadow-depth classes (G/L/C/M) classify entanglement
complexity: class G (Gaussian, $r_{\max}=2$) exhibits area-law
entanglement only; class L ($r_{\max}=3$) adds subleading logarithmic
corrections from the cubic shadow; class C ($r_{\max}=4$) adds a
contact correction from the quartic invariant; class M
($r_{\max}=\infty$) carries an infinite tower of corrections.

\medskip
\noindent\textbf{The entanglement programme (G11--G16).}\enspace
Six theorem targets extend the scalar result. G11: Ryu--Takayanagi
from the scalar projection of $\Theta_\cA$. G12: Koszulness as exact
quantum error correction; the Knill--Laflamme condition is the
Lagrangian isotropy condition from the complementarity pairing;
shadow depth equals the number of redundancy channels; the twelve
Koszulness characterisations translate into a twelve-fold code
dictionary; failure of Koszulness is failure of the code. G13:
replica trick from the MC equation. G14: modular flow from the
tangent complex. G15: QES as the shadow connection Ward identity.
G16: Page curve at $c=13$ from the shadow connection at the self-dual
point; the scrambling time is controlled by $\Delta_\cA$.


% ====================================================================
\section*{6.\quad The shadow obstruction tower}
% ====================================================================

\subsection*{6.1.\enspace The extension tower}

Can $\Theta_\cA$ be approximated by finitely many degrees, and what
controls the passage from one finite approximation to the next? The
total weight filtration answers both.

Define the \emph{total weight} $w(g,r,d):=2g-2+r+d$. The filtration
$F^N\mathfrak g^{\mathrm{amb}}_\cA:=\bigoplus_{w\ge N}(\mathfrak g^{\mathrm{amb}}_\cA)_{g,r,d}$
is exhaustive, separated, complete. $\mathfrak g^{\mathrm{amb}}_\cA$
is the \emph{ambient} (planted-forest-extended) modular convolution
algebra. The \emph{modular extension tower}:
\[
\cE_\cA(N):=\operatorname{MC}(\mathfrak g^{\mathrm{amb}}_\cA/F^{N+1}),\qquad N\ge 1,
\]
with restriction maps $\cE_\cA(N+1)\to\cE_\cA(N)$. A full lift is a
compatible point of the inverse system; the bar-intrinsic construction
provides one unconditionally. The shadow obstruction tower studies
$\cE_\cA(1)\leftarrow\cE_\cA(2)\leftarrow\cE_\cA(3)\leftarrow\cdots$.

\subsection*{6.2.\enspace Obstruction classes and the all-degree master equation}

Define $\Theta_\cA^{\le r}$ by projecting to total weight $\le r$. The
\emph{shadow obstruction tower} is the sequence
$\Theta_\cA^{\le 2},\Theta_\cA^{\le 3},\Theta_\cA^{\le 4},\ldots$ At
each truncation the MC equation fails by a remainder,
\[
D\Theta^{\le r}_\cA+\tfrac12[\Theta^{\le r}_\cA,\Theta^{\le r}_\cA]=-o_{r+1}(\cA),
\]
with obstruction class $o_{r+1}(\cA)\in Z^2(F^{r+1}/F^{r+2},d_{\le r})$.
Extension is possible iff $[o_{r+1}]=0$ in
$H^2(F^{r+1}/F^{r+2},d_{\le r})$. The bar-intrinsic construction
guarantees extension, so the obstruction class is always exact; the
shadow obstruction tower records the gauge choices at each stage.

The all-degree master equation at degree~$r$:
\[
\nabla_H(\operatorname{Sh}_r(\cA))+o^{(r)}(\cA)=0,\qquad r\ge 3,
\]
with $\nabla_H$ the Hodge covariant derivative. This equation
determines $\operatorname{Sh}_r$ uniquely modulo gauge if
$H^1(F^r/F^{r+1},d_{\le r-1})=0$.

\subsection*{6.3.\enspace Cubic gauge triviality and the canonical quartic}

\begin{quote}
\textbf{Cubic gauge triviality.}\enspace
When $H^1(F^3/F^4,d_2)=0$, the cubic MC term $\Theta^{[3]}_\cA$ is
gauge-trivial: there exists $\exp(\xi_3)$ with $\xi_3\in F^3$ such that
$\exp(\xi_3)\cdot\Theta^{\le 3}_\cA$ has vanishing weight-$3$
component.
\end{quote}

After this gauge transformation, the quartic obstruction class
$[\mathfrak Q(\cA)]\in H^2(F^4/F^5,d_2)$ is canonical: independent of
homotopy retraction data, cubic gauge, and MC gauge orbit
representative. This is the first characteristic class not polynomial
in $\kappa$. The mechanism is standard Postnikov: trivial $H^1$ at
level $n$ makes the $H^2$ class at level $n+1$ canonical.

The clutching law for the quartic follows from Mok's log FM
degeneration formula. For $C\leadsto C_0=C_1\cup_p C_2$,
\[
\mathfrak Q(\cA)|_{C_0}=\sum_{a+b=4}\mathfrak Q^{(a)}(\cA)|_{C_1}\cdot\mathfrak Q^{(b)}(\cA)|_{C_2}+\text{planted-forest correction}.
\]
The correction is genuinely logarithmic, with no DM counterpart.

The quartic contact invariant
$Q^{\mathrm{contact}}_{\mathrm{Vir}}=10/[c(5c+22)]\in H^2(F^4/F^5,d_2)$.
The denominator has physical meaning: $c=0$ is the logarithmic
threshold; $c=-22/5$ is the minimal model accumulation point
$\mathcal M(2,5)$. The numerator $10=\binom{5}{2}$ counts degree-four
Catalan channels.

\subsection*{6.4.\enspace The Hamilton--Jacobi generating function}

For a multi-generator algebra ($\mathcal W_3$ with generators $T$
(weight~$2$) and $W$ (weight~$3$)), the deformation space is
two-dimensional. The Riccati ODE lifts to a Hamilton--Jacobi PDE.
Write $U=\sum_{s\ge 3}\operatorname{Sh}_s(x_T,x_W)$; the master equation
becomes
\[
2E(U)+\tfrac12\lVert\nabla U\rVert^2_H=R(x_T,x_W),
\]
with Euler operator $E=x_T\partial_T+x_W\partial_W$, kinetic energy
$\lVert\nabla U\rVert^2_H=(2/c)(\partial_T U)^2+(3/c)(\partial_W U)^2$,
and $R$ a polynomial of degree $\le 4$ in
$(\kappa,\mathfrak C,\mathfrak Q)$. The gradient $dU\colon X\to T^*X$
embeds the deformation space as a Lagrangian submanifold of $T^*X$.
Whether the genus expansion admits a quantisation with $U$ as leading
semiclassical term is open: it would require a quantum Hamiltonian
whose WKB expansion has leading term $U$ and whose quantum corrections
reproduce the genus-$g$ shadow amplitudes.

On the $T$-line the conformal weight forces the transverse momentum to
vanish and the PDE collapses to the Riccati ODE. On the $W$-line the
cubic shadow $\operatorname{Sh}_3=2x_T^3+3x_Tx_W^2$ has nonzero normal
derivative $(\partial_T\operatorname{Sh}_3)|_{x_T=0}=3x_W^2$, and the
first inter-channel coupling correction enters at degree~$6$:
\[
\delta_6^{W\text{-line}}=\frac{576(5c+38)}{c^3(5c+22)^3}.
\]
The $T$-line is autonomous; the $W$-line is not. The asymmetry is
intrinsic to the weight difference between $T$ (weight~$2$) and $W$
(weight~$3$).

\subsection*{6.5.\enspace The four depth classes}

The \emph{shadow depth} $r_{\max}(\cA)$ is the smallest $r$ such that
$o_{r+1}(\cA)=0$ for all subsequent levels, or $\infty$ if the tower
never terminates. The standard landscape partitions into four depth
classes; no other values occur on any one-dimensional primary slice.
The single-line dichotomy: the shadow metric $Q_L(t)$ is quadratic in
$(\kappa,\alpha,S_4)$, and its factorisation type (perfect square or
irreducible) determines whether the tower terminates.

\begin{center}
\renewcommand{\arraystretch}{1.3}
\begin{tabular}{lcccll}
\textbf{Class} & \textbf{Symbol} & $r_{\max}$ & $\ell_3^{(0)},\ell_4^{(0)}$ & \textbf{Potential} & \textbf{Examples}\\
\hline
Gaussian & G & $2$ & $0,0$ & quadratic & Heisenberg $\cH_k$, free fermion $\psi$\\
Lie/tree & L & $3$ & $\neq 0,0$ & cubic & Affine $\widehat{\mathfrak{sl}}_N$, currents\\
Contact & C & $4$ & $0,\neq 0$ & quartic & $\beta\gamma$, $bc$ ghosts\\
Mixed & M & $\infty$ & $\neq 0,\neq 0$ & non-polynomial & Virasoro, $\mathcal W_N$, $\mathcal W_\infty$
\end{tabular}
\end{center}

\noindent
\textbf{Gaussian} ($r_{\max}=2$): transferred brackets $\ell_k^{(0)}$
vanish for $k\ge 3$ on bar cohomology, so the tower terminates at
$\kappa$. Two mechanisms: Heisenberg has no simple pole; the free
fermion has a simple pole but fermionic antisymmetry forces $S_3=0$.
Sewing is by Fredholm determinant (Heisenberg) or Pfaffian (fermion).
\textbf{Lie/tree} ($r_{\max}=3$): the OPE has a simple pole generating
a Lie bracket, but $\ell_4^{(0)}$ vanishes on bar cohomology.
\textbf{Contact} ($r_{\max}=4$): the simple-pole bracket vanishes on
bar cohomology, but the quaternary bracket survives; the first
nontrivial shadow is $\mathfrak Q(\cA)$. \textbf{Mixed}
($r_{\max}=\infty$): both cubic and quartic shadows are nontrivial,
and their Poisson bracket $\{\mathfrak C,\mathfrak Q\}_H$ generates a
nonvanishing quintic, then sextic, and so on.

\medskip
\noindent\textbf{Shadow depth is not Koszulness.}\enspace
All four depth classes contain chirally Koszul algebras. Heisenberg
(depth~$2$), affine at generic level (depth~$3$), $\beta\gamma$
(depth~$4$), and Virasoro (depth~$\infty$) are all chirally Koszul.
Shadow depth classifies the complexity of Koszul algebras, not
Koszulness itself.

The \emph{operadic complexity theorem}
(Theorem~\ref{thm:operadic-complexity-detailed}):
\[
r_{\max}(\cA)=A_\infty\text{-depth}(\cA)=L_\infty\text{-formality level}(\cA).
\]
$A_\infty$-depth is the largest $r$ with $m_r\neq 0$ in the minimal
$A_\infty$ model; $L_\infty$-formality level is the largest $r$ with a
nontrivial $r$-ary bracket on bar cohomology not removable by
$L_\infty$-isomorphism.

\subsection*{6.6.\enspace The genus--degree bigrading}

The MC equation decomposes along two orthogonal axes. The \emph{genus
axis} extends $\Theta_\cA$ from genus~$0$ to~$1$ to~$2$, controlled by
the genus spectral sequence. The \emph{degree axis} extends
$\Theta_\cA^{\le 2}\to\Theta_\cA^{\le 3}\to\Theta_\cA^{\le 4}$,
controlled by cyclic obstruction classes. At fixed degree~$2$ the
genus tower is governed by $\kappa$ alone \textsc{(uniform-weight)};
at fixed genus~$0$ each new shadow coefficient introduces independent
data. The axes couple through the visibility formula
$g_{\min}(S_r)=\lfloor r/2\rfloor+1$: $S_r$ first contributes at
genus $\lfloor r/2\rfloor+1$ and two new shadow coefficients enter at
each genus.

\subsection*{6.7.\enspace Genus-two shells}

Genus~$2$ is the first level where all three boundary strata interact:
\[
\Theta^{(2)}_\cA=\underbrace{\Theta^{(2)}_{\mathrm{loop}\circ\mathrm{loop}}}_{\text{iterated BV}}+\underbrace{\Theta^{(2)}_{\mathrm{sep}\circ\mathrm{loop}}}_{\text{clutching $\times$ BV}}+\underbrace{\Theta^{(2)}_{\mathrm{pf}}}_{\text{planted-forest}}.
\]
The iterated-BV term involves the figure-eight graph:
\[
\Theta^{(2)}_{\mathrm{loop}^2}=\tfrac12\operatorname{tr}_1\operatorname{tr}_2\bigl(P_\cA^{\otimes 2}\circ m^{(h)}_{0,4}\circ\Delta^{\otimes 2}\bigr).
\]
The separating-times-loop term:
\[
\Theta^{(2)}_{\mathrm{sep}\circ\mathrm{loop}}=\operatorname{tr}_{\mathrm{ns}}\bigl(P_\cA\circ\ell_2^{(0)}(\Theta^{(1,\bullet)},\Theta^{(1,\bullet)})\bigr).
\]
The planted-forest term is the genuinely logarithmic contribution,
present only when cubic or quartic shadows are nontrivial.

\medskip
\noindent
Family-by-family structure:

\begin{center}
\renewcommand{\arraystretch}{1.2}
\begin{tabular}{lcccc}
& $\Theta^{(2)}_{\mathrm{loop}^2}$ & $\Theta^{(2)}_{\mathrm{sep}\circ\mathrm{loop}}$ & $\Theta^{(2)}_{\mathrm{pf}}$ & \textbf{Depth}\\
\hline
Heisenberg $\cH_k$ & $\checkmark$ & --- & --- & G\\
Affine $\widehat{\fg}_k$ & $\checkmark$ & $\checkmark$ & --- & L\\
$\beta\gamma$ system & $\checkmark$ & --- & --- & C\\
Virasoro $\mathrm{Vir}_c$ & $\checkmark$ & $\checkmark$ & $\checkmark$ & M\\
$\mathcal W_N$ & $\checkmark$ & $\checkmark$ & $\checkmark$ & M
\end{tabular}
\end{center}

Heisenberg: only iterated BV. Affine: adds separating-times-loop.
Virasoro and $\mathcal W_N$: all three active, with planted-forest
carrying the quartic resonance. $\beta\gamma$: only iterated BV
survives; the separating term vanishes because the Lie bracket is
trivial on bar cohomology, and the planted-forest vanishes because
the quartic contact invariant $\mu_{\beta\gamma}=0$ (rank-one abelian
rigidity). Contact termination via the iterated BV shell, not
Gaussian collapse.

\subsection*{6.8.\enspace The genus loop operator}

The \emph{genus loop operator}
$\Lambda_P\colon\operatorname{Sym}^r(\bar\cA^\vee)_{\mathrm{cyc}}\to\operatorname{Sym}^{r-2}(\bar\cA^\vee)_{\mathrm{cyc}}$
contracts two legs with the propagator:
\[
\Lambda_P(\phi)(a_1,\ldots,a_{r-2})=\sum_i\phi(a_1,\ldots,a_{r-2},e_i,e^i).
\]
This raises genus by~$1$ and lowers degree by~$2$: the shadow-level
incarnation of $\Delta_{\mathrm{ns}}$.

\medskip
\noindent\textbf{Explicit Virasoro computation.}\enspace
On the Virasoro primary line, the conformal-weight-$2$ bar cohomology
is one-dimensional, spanned by the class $x$ dual to $T$. The Killing
form is $H=(c/2)x^2$, the propagator $P=2/c$, and the quartic contact
invariant is $Q^{\mathrm{contact}}_{\mathrm{Vir}}\cdot x^4$ with
$Q^{\mathrm{contact}}_{\mathrm{Vir}}=10/[c(5c+22)]$. Applying
$\Lambda_P$:
\[
\Lambda_P\!\left(\frac{10}{c(5c+22)}\cdot x^4\right)=\binom{4}{2}\cdot\frac{2}{c}\cdot\frac{10}{c(5c+22)}\cdot x^2=\frac{120}{c^2(5c+22)}x^2,
\]
the genus-one Hessian correction. The loop ratio
\[
\rho^{(1)}_{\mathrm{Vir}}=\frac{\Lambda_P(Q\cdot x^4)}{\kappa\cdot x^2}=\frac{120/[c^2(5c+22)]}{c/2}=\frac{240}{c^3(5c+22)}
\]
is not determined by $\kappa=c/2$: it detects the quartic structure
through its genus-one projection. Asymptotically
$\rho^{(1)}\sim 48/c^4$ for $c\to\infty$.

\subsection*{6.9.\enspace Quintic forcing and infinite towers}

The quintic obstruction determines whether the tower terminates at
degree~$4$.

\medskip
\noindent\textbf{Virasoro.}\enspace
Both $\mathfrak C_{\mathrm{Vir}}$ and $Q^{\mathrm{contact}}_{\mathrm{Vir}}$
are nonzero. Their Poisson bracket
\[
o^{(5)}_{\mathrm{Vir}}=\{\mathfrak C_{\mathrm{Vir}},Q^{\mathrm{contact}}_{\mathrm{Vir}}\}_H=\frac{480}{c^2(5c+22)}x^5\neq 0.
\]
For all $c\neq 0,-22/5$, the Virasoro shadow tower does not terminate;
$r_{\max}(\mathrm{Vir}_c)=\infty$. Same for every $\mathcal W_N$.

\medskip
\noindent\textbf{Heisenberg.}\enspace
$\mathfrak C_{\cH}=0$, $Q^{\mathrm{contact}}_{\cH}=0$, hence $o^{(5)}=0$
and all higher obstructions vanish. $r_{\max}=2$.

\medskip
\noindent\textbf{Affine.}\enspace
$\mathfrak C_{\widehat\fg}\neq 0$ but $Q^{\mathrm{contact}}_{\widehat\fg}=0$
(quaternary bracket vanishes on bar cohomology by the PBW theorem).
$o^{(5)}=0$; $r_{\max}=3$.

\medskip
\noindent\textbf{$\beta\gamma$.}\enspace
$\mathfrak C_{\beta\gamma}=0$ (Lie bracket trivial on bar cohomology)
but $Q^{\mathrm{contact}}_{\beta\gamma}\neq 0$. $o^{(5)}=0$;
$r_{\max}=4$.

\begin{center}
\renewcommand{\arraystretch}{1.15}
\begin{tabular}{lcccc}
& $\mathfrak C$ & $\mathfrak Q$ & $o^{(5)}=\{\mathfrak C,\mathfrak Q\}_H$ & $r_{\max}$\\
\hline
Heisenberg & $0$ & $0$ & $0$ & $2$\\
Affine & $\neq 0$ & $0$ & $0$ & $3$\\
$\beta\gamma$ & $0$ & $\neq 0$ & $0$ & $4$\\
Virasoro & $\neq 0$ & $\neq 0$ & $\neq 0$ & $\infty$
\end{tabular}
\end{center}

Finite termination requires $\mathfrak C=0$ or $\mathfrak Q=0$; when
both are nonzero, their bracket generates the quintic, and the
induction continues without bound.

\medskip
\noindent\textbf{The shadow growth rate.}\enspace
For class~M algebras, the Taylor coefficients $S_r$ of $\sqrt{Q_L}$
grow as $S_r\sim A\cdot\rho(\cA)^r\cdot r^{-5/2}\cos(r\theta+\varphi)$
with shadow growth rate $\rho(\cA)=\sqrt{9\alpha^2+2\Delta}/(2|\kappa|)$,
a continuous invariant. Self-dual Virasoro ($c=13$) has
$\rho\approx 0.467$. The critical cubic $5c^3+22c^2-180c-872=0$ has
root $c^*\approx 6.125$: the tower converges for $c>c^*$ and diverges
for $c<c^*$.

\medskip
\noindent\textbf{The shadow partition function.}\enspace
The double sum $Z^{\mathrm{sh}}_\cA=\sum_{g,r}\hbar^{2g-2}Z_g^{(r)}$
converges absolutely. Genus decay at rate $1/(2\pi)^{2g}$
(Faber--Pandharipande); degree decay at rate $\rho^r r^{-5/2}$
(Darboux analysis). The string partition function diverges as
$(2g-2)!$. For Virasoro at $c=26$, $\rho\approx 0.234$ and
$Z^{\mathrm{sh}}$ has positive radius of convergence in both $\hbar$
and degree.

\subsection*{6.10.\enspace Independent sum factorisation}

For a direct sum $L=L_1\oplus L_2$ with vanishing mixed OPE, all
shadow invariants separate:
\begin{align*}
\kappa(L_1\oplus L_2)&=\kappa(L_1)+\kappa(L_2),\\
T^{\mathrm{br}}(L_1\oplus L_2)&=T^{\mathrm{br}}(L_1)\oplus T^{\mathrm{br}}(L_2),\\
\Delta_{L_1\oplus L_2}(t)&=\Delta_{L_1}(t)\cdot\Delta_{L_2}(t),\\
R_4^{\mathrm{mod}}(L_1\oplus L_2)&=R_4^{\mathrm{mod}}(L_1)+R_4^{\mathrm{mod}}(L_2).
\end{align*}
The full MC element factorises:
$\Theta_{L_1\oplus L_2}=\Theta_{L_1}\hat\otimes 1+1\hat\otimes\Theta_{L_2}$,
and $Z_{L_1\oplus L_2}=Z_{L_1}\cdot Z_{L_2}$.

Every algebra in the standard landscape is either primitive
(Heisenberg, Virasoro, $\beta\gamma$) or built from primitive pieces
(affine = current algebra over simple Lie, lattice = exponential
extension of Heisenberg, $\mathcal W_N$ = quantum Drinfeld--Sokolov
reduction of affine). Drinfeld--Sokolov reduction is the
\emph{shadow-depth escalator}
(Principle~\ref{princ:shadow-depth-escalator}): it sends current
algebras (class L, $r_{\max}=3$, finite towers) to $\mathcal W$-algebras
(class M, $r_{\max}=\infty$, infinite towers). The nonlinearity of
the Dirac bracket on the Slodowy slice creates the quartic and all
higher shadows.

\subsection*{6.11.\enspace The representation-theoretic perspective}

The categorical logarithm $\eta_{ij}=d\log(z_i-z_j)$ of \S1 is the
integral kernel of a tautology: $D_\cA^2=0$ on the modular bar
complex. Three structure theorems emerge. \emph{Algebraicity}: on
each primary slice $L$, $H(t)^2=t^4Q_L(t)$ with $Q_L$ quadratic in
$\kappa$, $\alpha$, $S_4$. \emph{Formality}: the genus-$0$ shadow
obstruction tower is the $L_\infty$ formality obstruction tower; the
four depth classes G/L/C/M are the corresponding formality types,
while positive-genus corrections form a separate quantum layer.
\emph{Complementarity}: for a Koszul pair $(\cA,\cA^!)$, the summands
$\mathbf Q_g(\cA)\oplus\mathbf Q_g(\cA^!)$ are Lagrangian (Theorem~C).

The discriminant $\Delta=8\kappa S_4$ governs the single-line
dichotomy: $\Delta=0$ forces tower termination
($L_\infty$-formality); $\Delta\neq 0$ forces an infinite tower
(intrinsic non-formality). The contact class C ($\beta\gamma$)
escapes via stratum separation. No invariant of the classical bar
complex of a graded algebra sees~$\Delta$: it is visible only on
curves of positive genus, where the Arnold relation acquires a defect.

Classical Koszul duality (Priddy, Beilinson--Ginzburg--Soergel)
embeds into the genus~$0$, degree~$2$, $\Delta=0$ stratum via
formal-disk restriction (the formal, class-G case); the
$E_1/E_\infty$ distinction and Fulton--MacPherson geometry are
already absent over a bare point.
The $R$-matrix of quantum groups arises as the genus-$0$ shadow
$\operatorname{Sh}_{0,2}(\Theta_{\widehat\fg_k})$ (\S7.7). At the
critical level $k=-h^\vee$, the scalar class vanishes ($\kappa=0$)
and bar cohomology recovers the Feigin--Frenkel center
$\operatorname{Fun}(\operatorname{Op}_{{}^L\fg})$; the Langlands
programme lives on this orthogonal axis, not in the MC element but
in the cohomology that emerges on the uncurved scalar locus.


% ====================================================================
\section*{9.\quad Scalar saturation and the Koszulness programme}
% ====================================================================

The bar construction works for all chiral algebras. It works best,
the cobar inverts the bar, the spectral sequence collapses, the
resolution is linear, precisely on the Koszul locus. Characterizing
this locus is the central foundational question introduced in~\S3.3
and anticipated throughout sections~4 and~5. Two complementary
results: scalar saturation (the MC element is determined by a single
scalar for all standard families) and the Koszulness meta-theorem
(twelve characterisations, ten unconditionally equivalent, one
conditional on perfectness, one one-directional).

\subsection*{9.1.\enspace Scalar saturation}

For a one-parameter algebraic family $\cA_k$ with rational OPE
coefficients (structure constants rational functions of~$k$), the
universal MC element satisfies
\[
\Theta_{\cA_k}=\kappa(\cA_k)\cdot\eta_{\mathrm{univ}}\otimes\Lambda_{\mathrm{univ}}+(\text{terms in }F^3\mathfrak g^{\mathrm{mod}}_{\cA_k}).
\]
\emph{Scalar saturation universality}: higher-order terms vanish. For
every algebraic family with rational OPE coefficients, $\Theta_{\cA_k}$
is determined by $\kappa(\cA_k)$ modulo terms in the degree-$\ge 3$
filtration.

\medskip
\noindent\textbf{Step~1: Whitehead reduction.}\enspace
For a vertex algebra with reductive symmetry at generic level, the
primitive cyclic cohomology
$H^2_{\mathrm{cyc,prim}}(\mathfrak g^{\mathrm{mod}}/F^3,d_2)$ is
computed by a Whitehead-type spectral sequence with
$E_2^{p,q}=H^p(\mathfrak g;\operatorname{Sym}^q V)$. For reductive
$\mathfrak g$ and generic level, $H^2(\mathfrak g;\operatorname{Sym}^q V)=0$
by the Whitehead lemma, and the degree-$2$ MC class lifts uniquely.

\medskip
\noindent\textbf{Step~2: Algebraic semicontinuity.}\enspace
The vanishing locus
\[
E=\{k\in\mathbb C:H^2_{\mathrm{cyc,prim}}(\cA_k)\neq 0\}
\]
is Zariski-closed by upper semicontinuity. At $k=\infty$ the algebra
is Heisenberg and $H^2_{\mathrm{cyc,prim}}=0$, so $E$ is a proper
Zariski-closed subset. Explicit computation verifies $E=\varnothing$
for $\mathfrak{sl}_N$ ($N=2,\ldots,10$): no exceptional level. The
result covers the entire standard Lie-theoretic landscape at all
non-critical levels, including admissible.

\medskip
\noindent\textbf{Scope.}\enspace
Scalar saturation applies to every family in the standard landscape:
affine Kac--Moody (all types, all non-critical levels), Virasoro (all
$c$), $\mathcal W_N$ (all $N$, all non-degenerate levels),
$\beta\gamma$ and lattice algebras (all parameters). The residual
conjecture, that $E=\varnothing$ for every positive-energy chiral
algebra with rational OPE coefficients, is open only for
non-algebraic-family constructions.

\subsection*{9.2.\enspace The Koszulness meta-theorem}

The Koszul condition admits twelve formulations, spanning homological
algebra (PBW collapse, Ext vanishing), homotopy theory ($A_\infty$
formality, $E_2$ formality), operadic algebra
(Barr--Beck--Lurie, FM acyclicity), modular geometry (FH
concentration, tropical Koszulness), and symplectic topology
(Lagrangian criterion). Ten are unconditionally equivalent; one (the
Lagrangian criterion) is conditional on perfectness/nondegeneracy;
one ($\mathcal D$-module purity) is one-directional (forward direction
proved, converse open).

\begin{center}
\emph{For $\cA$ positive-energy and finitely strongly generated, the
following are equivalent:}
\end{center}

\begin{enumerate}[label=\textup{(\roman*)},nosep]
\item $\cA$ is chirally Koszul.
\item The PBW spectral sequence on $\barB^{(g)}(\cA)$ collapses at $E_2$ for every $g$.
\item The minimal $A_\infty$-model of $H^\bullet(\barB(\cA))$ is formal: $m_n=0$ for $n\ge 3$.
\item $\operatorname{Ext}^{p,q}_\cA(\omega_X,\omega_X)=0$ for $p\neq q$.
\item The bar-cobar counit $\Omega(\barB(\cA))\to\cA$ is a quasi-isomorphism.
\item The Barr--Beck--Lurie comparison for $\barB\dashv\Omega$ is an equivalence.
\item $\int_X\cA$ is concentrated in degree~$0$ for every smooth proper $X$.
\item $\operatorname{ChirHoch}^*(\cA)$ vanishes outside degrees $\{0,1,2\}$.
\item The Kac--Shapovalov determinant $\det G_h\neq 0$ in the bar-relevant range.
\item The restriction $i_S^!\barB_n(\cA)$ to every FM boundary stratum $S$ is acyclic outside degree~$0$.
\end{enumerate}

\noindent
Under the additional hypothesis that the ambient tangent complex is
perfect and nondegenerate:

\begin{enumerate}[label=\textup{(\roman*)},nosep,start=11]
\item $\mathcal M_\cA$ and $\mathcal M_{\cA^!}$ are transverse Lagrangians in the $(-1)$-shifted symplectic deformation space $\mathcal M_{\mathrm{comp}}$.
\end{enumerate}

\noindent
Finally (xii)$\Rightarrow$(x), where:

\begin{enumerate}[label=\textup{(\roman*)},nosep,start=12]
\item Each $\barB_n(\cA)$ is pure of weight~$n$ as a mixed Hodge module, with characteristic variety aligned to FM boundary strata ($\mathcal D$-module purity).
\end{enumerate}

\noindent The converse of (xii) is open.

\medskip
\noindent\textbf{The proof circuit.}\enspace
The core equivalence chain is
(i)$\Leftrightarrow$(ii)$\Leftrightarrow$(iii)$\Leftrightarrow$(v)$\Leftrightarrow$(viii),
proved by PBW concentration $\to$ bar-cobar inversion $\to$ Hochschild
polynomial growth $\to$ Koszul-complex acyclicity, with the $A_\infty$
formality branch via HPL transfer and Keller classicality. From this
core: (iv) by Ext spectral sequence, (vi) by conservativity with
log-FM totalization, (vii) by $\barB=\int_X$, (ix) by non-degeneracy
forcing PBW injectivity, (x) by stratum-by-stratum PBW concentration.
The Lagrangian criterion (xi) uses PTVV derived intersection theory
on the perfectness hypothesis. The $\mathcal D$-module purity
implication (xii)$\Rightarrow$(x) uses weight filtration on the bar
complex; the converse would require concentration to force purity.
The cycle closes via (viii)$\Rightarrow$(i): Koszul-complex acyclicity
implies the defining characterisation.

\medskip
\noindent\textbf{Beyond the meta-theorem.}\enspace
Six additional characterisations are proved in the text:

\begin{enumerate}[label=\textup{(\alph*)},nosep]
\item \emph{Shadow-formality at low degree.} The shadow obstruction tower $\Theta^{\le r}_\cA$ equals the $L_\infty$-formality obstruction tower of $\mathfrak g^{\mathrm{mod}}_\cA$ at degrees $r=2,3,4$
(Proposition~\ref{prop:shadow-formality-low-degree}). Full equivalence to~(iii) requires all degrees.

\item \emph{$E_2$-formality of chiral Hochschild cohomology.} Koszulness implies $\operatorname{ChirHoch}^*(\cA)$ is formal as an $E_2$-algebra (Deligne conjecture applied to the bar-cobar resolution).

\item \emph{Genus-$0$ curve independence.} Koszulness on $\mathbb P^1$ implies Koszulness on every smooth proper curve (Proposition~\ref{prop:genus0-curve-independence}). Proved independently by factorization descent.

\item \emph{PBW universality} (sufficient). Every freely strongly generated vertex algebra is chirally Koszul (Proposition~\ref{prop:pbw-universality}). Applies to universal algebras: $V_k(\mathfrak g)$ is Koszul at every level including critical and admissible. For simple quotients $L_k(\fg)$ at admissible levels, status is rank-dependent: $L_k(\mathfrak{sl}_2)$ is Koszul at all admissible levels, but at $\operatorname{rk}(\fg)\ge 2$ with denominator $q\ge 3$, Koszulness fails via a $\operatorname{rk}(\fg)$-dimensional bar obstruction from the abelian Cartan
(Conjecture~\ref{conj:admissible-koszul-rank-obstruction}). PBW universality is one-directional: not every Koszul algebra is freely strongly generated.

\item \emph{Tropical Koszulness.} $\cA$ is chirally Koszul iff $\barB^{\mathrm{trop}}(\cA)$ is acyclic. Proved by the weight spectral sequence on the real blowup.

\item \emph{Bifunctor obstruction decomposition} (one-directional). Koszulness implies the RNW one-slot obstruction decomposes into independent bar-side and cobar-side components
(Theorem~\ref{thm:bifunctor-obstruction-decomposition}). The converse is open.
\end{enumerate}

\medskip
\noindent\textbf{Status summary.}
\begin{center}
\renewcommand{\arraystretch}{1.15}
\begin{tabular}{clll}
\multicolumn{4}{l}{\emph{Meta-theorem (Theorem~\ref{thm:koszul-equivalences-meta}):}}\\[1pt]
& (i)--(x) & \textbf{10 unconditional equivalences} & proved\\
& (xi) Lagrangian & conditional on perfectness & proved (unconditional for standard landscape)\\
& (xii) $\mathcal D$-module purity & (xii)$\Rightarrow$(x) only & converse open\\[3pt]
\multicolumn{4}{l}{\emph{Additional characterisations:}}\\[1pt]
& (a) Shadow-formality & equivalence at degrees $2,3,4$ & proved\\
& (b) $E_2$-formality & consequence of Koszulness & proved\\
& (c) Curve independence & independent theorem & proved\\
& (d) PBW universality & sufficient condition & proved\\
& (e) Tropical Koszulness & equivalence & proved\\
& (f) Bifunctor decomposition & one-directional & converse open
\end{tabular}
\end{center}

\subsection*{9.3.\enspace Shadow depth and operadic complexity}

The four depth classes G/L/C/M admit refinement: in general
$r_{\max}\in\{2,3,4,5,\ldots\}\cup\{\infty\}$, with class $\mathbf F_r$
(tower terminates at degree~$r$) structurally permitted for all
$r\ge 2$; $\mathbf G=\mathbf F_2$, $\mathbf L=\mathbf F_3$,
$\mathbf C=\mathbf F_4$ name the standard-landscape representatives.

\medskip
\noindent\textbf{The operadic complexity theorem}
(Theorem~\ref{thm:operadic-complexity-detailed}):
\[
r_{\max}(\cA)=A_\infty\text{-depth}(\cA)=L_\infty\text{-formality level}(\mathfrak g^{\mathrm{mod}}_\cA).
\]
$A_\infty$-depth is the minimal $r$ with $m_k=0$ for $k>r$ on
$H^\bullet(\barB(\cA))$; $L_\infty$-formality level is the minimal $r$
such that the $L_\infty$-structure on $\mathfrak g^{\mathrm{mod}}_\cA$
is formal through degree~$r$. Three consequences: the shadow
obstruction tower \emph{is} the $L_\infty$-formality obstruction
tower; the depth classification becomes a structural invariant of the
operadic type of $\cA$; the tower of a Gaussian algebra is a quadratic
form, of a Lie algebra a Chevalley--Eilenberg complex, of a contact
algebra a cyclic $A_\infty$-structure, of a mixed algebra an infinite
$L_\infty$-tower with no polynomial truncation.

The programme extends to non-principal $\mathcal W$-algebras: the
Bershadsky--Polyakov algebra $\mathrm{BP}_k$ is chirally Koszul at generic level ($k \neq -3$) with
dual $\mathrm{BP}_{-k-6}$ and Koszul conductor
$K_{\mathrm{BP}}=196$. Nilpotent transport from hook-type seeds
covers all orbits in $\mathfrak{sl}_2$ through $\mathfrak{sl}_8$.

\subsection*{9.4.\enspace The PBW spectral sequence}

The computational engine for bar cohomology within fixed genus is the
PBW (Poincar\'e--Birkhoff--Witt) spectral sequence. Filter
$\barB^{(g)}(\cA)$ by \emph{conformal weight}: assign each normally
ordered monomial its total weight (sum of constituent field weights).
The associated graded is the bar complex of a free-field algebra,
reducing to symmetric-function combinatorics.

The $E_1$ page at genus~$g$:
\[
E_1^{p,q}(g)=E_1^{p,q}(g{=}0)\oplus\cE^{p,q}_g,
\]
with the first summand the genus-$0$ bar cohomology determined by
Koszulness of the associated graded and $\cE^{p,q}_g$ the genus-$g$
enrichment from the Hodge bundle, arising from the regular Eisenstein
series coupled to the curvature. \emph{Concentration}:
$E_\infty^{p,q}(g)=0$ for $q\neq 0$ at all $g\ge 1$, for all standard
families. PBW degeneration: the spectral sequence collapses at $E_2$
and bar cohomology concentrates on the diagonal.

The PBW spectral sequence is distinct from the genus spectral sequence
of \S5.5. PBW uses conformal weight within fixed genus to compute bar
cohomology groups; the genus spectral sequence uses the genus
filtration on the deformation algebra to control the MC lift across
genera. One computes \emph{what}; the other determines \emph{whether}.

\subsection*{9.5.\enspace Chiral Hochschild cohomology (Theorem~H)}

The chiral Hochschild complex $\operatorname{ChirHoch}^\bullet(\cA)$
is the derived endomorphism algebra of the identity functor on
$\cA$-modules, computed by the bar construction with bimodule
coefficients. At genus~$0$,
\[
\operatorname{ChirHoch}^\bullet(\cA)\simeq R\Hom_{\cA\text{-bimod}}(\cA,\cA)\simeq\operatorname{Ext}^\bullet_{\cA^e}(\cA,\cA),
\]
where $\cA^e=\cA\otimes\cA^{\mathrm{op}}$ is the enveloping chiral
algebra.

For $\cA=\mathcal W^k(\mathfrak g,f_{\mathrm{prin}})$ a principal
$\mathcal W$-algebra at generic level, the chiral Hochschild cohomology
is concentrated in degrees $\{0,1,2\}$:
\[
\operatorname{ChirHoch}^n(\mathcal W^k(\mathfrak g))=0\text{ for }n>2,\qquad P(t)=1+t^2.
\]
The deformation class in $\operatorname{ChirHoch}^2=\mathbb C$ is the
level deformation $k\mapsto k+\epsilon$. The concentration bound is
the de Rham amplitude on a curve ($\dim X=1$); the continuous
Lie algebra cohomology
$H^*_{\mathrm{cont}}(\operatorname{Lie}(\mathcal W),\operatorname{Lie}(\mathcal W)_0;\mathbb C)=\mathbb C[\Theta_1,\ldots,\Theta_r]$
(Gelfand--Fuchs, $\deg\Theta_i=m_i+1$) is a different, unbounded
invariant. For $\mathfrak g=\mathfrak{sl}_2$ (Virasoro):
$P(t)=1+t^2$. For $\mathfrak g=\mathfrak{sl}_3$ ($\mathcal W_3$):
$P(t)=1+t^2$. Total dimension $\le 4$, never $\mathbb C[\Theta]$.

On the Koszul locus, $\operatorname{ChirHoch}^\bullet(\cA)$ carries an
$E_2$-algebra structure that is formal. The Connes periodicity
operator $S\colon\operatorname{ChirHoch}^n\to\operatorname{ChirHoch}^{n-2}$
intertwines Koszul duality: $S\circ\operatorname{KD}=\operatorname{KD}\circ S$.
At genus~$0$, the chiral Hochschild complex identifies with the bulk
chiral homology complex of Volume~II. The bulk-boundary-line triangle
is the genus-$0$ shadow of this identification:
$\cA_{\mathrm{bulk}}\simeq Z_{\mathrm{der}}(\cB_\partial)\simeq\operatorname{ChirHoch}^\bullet(\cA^!)$.

\subsection*{9.6.\enspace The shadow algebra}

The \emph{shadow algebra} $\cA^{\mathrm{sh}}$ is the cohomology of the
modular cyclic deformation complex as a graded-commutative ring,
\[
\cA^{\mathrm{sh}}:=H_\bullet(\Defcyc^{\mathrm{mod}}(\cA))=\bigoplus_{r\ge 2,g\ge 0}\cA^{\mathrm{sh}}_{r,g}.
\]
The single MC element $\Theta_\cA$ decomposes by degree and genus.
Named shadows are projections:
\[
\kappa(\cA)=\cA^{\mathrm{sh}}_{2,0},\quad\mathfrak C(\cA)=\cA^{\mathrm{sh}}_{3,0},\quad\mathfrak Q(\cA)=\cA^{\mathrm{sh}}_{4,0},\quad\operatorname{Sh}_r(\cA)=\cA^{\mathrm{sh}}_{r,0}\text{ for }r\ge 5.
\]
The shadow algebra organises the entire hierarchy; it is a homotopy
invariant, computed by the Weil homomorphism of the Chern--Weil
dictionary (\S5.4).

The all-degree master equation in the shadow algebra:
$\nabla_H(\operatorname{Sh}_r)+o^{(r)}=0$ for all $r\ge 3$. The
obstruction class $o^{(r)}$ lies in $\cA^{\mathrm{sh}}_{r,0}$ and
measures the failure of $\Theta^{\le r-1}_\cA$ to extend to degree~$r$.

\subsection*{9.7.\enspace BV/BRST identification at genus~$0$}

At genus~$0$, the bar differential of a chiral algebra on $\mathbb C$
is the BRST differential of the associated holomorphic twist,
$(d_{\barB})_{g=0}=Q_{\mathrm{BRST}}$. The bar complex
$\barB^{(0)}_{\mathbb C}(\cA)$ is the BRST complex of the 2d chiral
algebra, and bar-cobar inversion
$\Omega(\barB(\cA))\xrightarrow{\sim}\cA$ (Theorem~B) says the BRST
cohomology processed through the cobar functor reconstructs the
original algebra.

At genus~$1$, the identification acquires curvature:
$d^2_{\mathrm{fib}}=\kappa(\cA)\cdot\omega_1$ is the one-loop anomaly,
and the curved bar complex is the curved BRST complex of the theory
on an elliptic curve. The complementarity
$Q_1(\cA)\oplus Q_1(\cA^!)=H^*(\overline{\cM}_{1,1},\cZ_\cA)$ is the
genus-$1$ anomaly cancellation condition.

At genus $g\ge 2$, the identification with Costello--Gwilliam
requires renormalisation: higher-loop integrals over
$\overline{\cM}_{g,n}$ need regularisation beyond the algebraic
bar-cobar framework, and the correct receptacle is the coderived
category.

For the Heisenberg algebra, the identification is proved at all genera
by four independent paths: zeta-regularisation, Selberg zeta function,
Grothendieck--Riemann--Roch, direct bar computation. The shadow
obstruction tower equals the $L_\infty$-formality obstruction tower
at all degrees.

\medskip
\noindent\textbf{Genus-two shell activation as depth diagnostic.}\enspace
The genus-two shell profile independently confirms the depth
classification:

\begin{center}
\renewcommand{\arraystretch}{1.15}
\begin{tabular}{lccc}
& $\Theta^{(2)}_{\mathrm{loop}^2}$ & $\Theta^{(2)}_{\mathrm{sep}\circ\mathrm{loop}}$ & $\Theta^{(2)}_{\mathrm{pf}}$\\
\hline
\textbf{G} (Heisenberg, lattice) & $\checkmark$ & --- & ---\\
\textbf{L} (affine $\widehat{\mathfrak g}_k$) & $\checkmark$ & $\checkmark$ & ---\\
\textbf{C} ($\beta\gamma$) & $\checkmark$ & --- & ---\\
\textbf{M} (Virasoro, $\mathcal W_N$) & $\checkmark$ & $\checkmark$ & $\checkmark$
\end{tabular}
\end{center}

\noindent
G and C activate only the loop$\circ$loop shell; L adds sep$\circ$loop;
only M activates all three including planted-forest. The
planted-forest shell is the diagnostic: its presence forces infinite
shadow depth.

\medskip
\noindent\textbf{Structural refinements.}\enspace
The three bar complexes on a chiral algebra (ordered tensor $T^c$,
symmetric $\operatorname{Sym}^c$, Chevalley--Eilenberg
$\operatorname{Lie}^c$) sit in a chain
$\operatorname{Lie}^c\hookrightarrow\operatorname{Sym}^c\hookrightarrow T^c$
with distinct coalgebra structures at each level
(Theorem~\ref{thm:three-bar-complexes}). The averaging map
$\mathrm{av}\colon T^c\to\operatorname{Sym}^c$ from ordered to
symmetric is surjective but not split: the kernel at degree~$3$
contains the Drinfeld associator $\Phi_{\mathrm{KZ}}(\cA)$, and the
fibre of splittings is a torsor for $\mathrm{GRT}_1$. The $E_1$
ordered bar complex is therefore the natural primitive object; the
modular/symmetric story is its $\Sigma_n$-coinvariant image
(Theorem~\ref{thm:e1-primacy}). The genus-$3$ cross-channel
correction for $\mathcal W_3$ has been computed by summing over all
$42$ stable graphs of $\overline{\cM}_{3,0}$
(Proposition~\ref{prop:w3-genus3-cross-channel}): the cross-channel
tower accelerates with genus and dominates the scalar part
$\kappa\cdot\lambda_g^{\mathrm{FP}}$ by genus~$4$.

On the open/closed side, the Koszul dual cooperad
$\mathsf{SC}^{\mathrm{ch,top},!}$ decomposes into three sectors:
closed ($\dim=(n-1)!$, Lie cooperad), open ($\dim=m!$, Ass cooperad),
and mixed ($\dim=(k-1)!\binom{k+m}{m}$). The mixed sector
encodes the bulk-to-boundary module structure by which the chiral
derived centre acts on the boundary algebra.

The BV/BRST${}={}$bar identification at higher genus
(Conjecture~\ref{conj:master-bv-brst}): the chain-level obstruction
(harmonic propagator correction) is resolved for classes G, L, and C,
for G and L by the Heisenberg sewing theorem and its extension to
affine algebras, for C by contact decoupling. For class M (Virasoro,
$\mathcal W_N$) the identification holds on bar cohomology but fails
at the chain level: the harmonic propagator produces a nonvanishing
quartic obstruction, so a coderived reformulation is needed.
